% SSEE-V3.6 — Paper 1: Zero-Parameter Framework
% Compile: pdflatex (×2)
\documentclass[12pt,a4paper]{article}

\usepackage[utf8]{inputenc}
\usepackage{amsmath,amssymb,amsfonts}
\usepackage{graphicx}
\usepackage{booktabs}
\usepackage{hyperref}
\usepackage{xcolor}
\usepackage{geometry}
\usepackage{natbib}
\usepackage{caption}
\usepackage{float}
\usepackage{array}
\usepackage{enumitem}

\geometry{margin=2.5cm}
\hypersetup{colorlinks=true,linkcolor=blue,citecolor=blue,urlcolor=blue}

\newcommand{\SSEE}{\ifmmode\mathrm{SSEE}\else\textsc{ssee}\fi}
\newcommand{\LCDM}{\ifmmode\Lambda\mathrm{CDM}\else$\Lambda$CDM\fi}
\newcommand{\KALz}{\mathrm{KAL}_0}
\newcommand{\wz}{w_0}
\newcommand{\wa}{w_a}
\newcommand{\OmDE}{\Omega_{\mathrm{DE}}}
\newcommand{\Migimf}{M^{\mathrm{IGIMF}}_{\mathrm{bar}}}
\newcommand{\Mobs}{M^{\mathrm{obs}}_{\mathrm{dyn}}}

\title{\textbf{A Zero-Parameter Framework for Cosmological Dynamics\\
and Galaxy Cluster Mass Discrepancies via\\
Structural Self-Energy Expansion (SSEE-V3.6)}}

\author{Mike Edison Almeida Vallejo%
\thanks{Independent researcher.
E-mail: \href{mailto:mike.almeida1721@gmail.com}{mike.almeida1721@gmail.com}.
ORCID: \href{https://orcid.org/0009-0008-2195-7836}{0009-0008-2195-7836}.}%
}
\date{April 2026}

\begin{document}

\maketitle

\begin{abstract}
We present the Structural Self-Energy Expansion (\SSEE-V3.6), a minimalist theoretical
framework that unifies cosmological expansion and galaxy-cluster dynamics using only two
dimensionless geometric constants — $\pi$ and the golden ratio $\varphi$ — without
introducing free parameters or exotic dark-matter particles.
From these two axioms, a closed set of seven structural registers is derived algebraically,
culminating in the dark-energy equation of state
\begin{equation*}
  w_0 = -\frac{3\varphi+\pi}{2(\varphi+\pi)} \approx -0.840,
  \qquad
  w_a = -\frac{2\varphi+\pi}{2(\varphi+\pi)} \approx -0.670.
\end{equation*}
No fitting procedure is involved: $w_0$ and $w_a$ are uniquely fixed by the geometry of
$\varphi$ and $\pi$.
The asymptotic structural viscosity factor $\KALz = \beta + \pi \approx 5.5214$
natively reproduces the observed cosmic dark-to-baryonic matter density ratio.
At galactic scales, \SSEE\ derives an acceleration threshold
$a_0 = c\,H_0/\KALz \approx 1.186\times10^{-10}$\,m\,s$^{-2}$, consistent with
the empirical MOND scale.
Incorporating a light-neutrino mass sum $\sum m_\nu = 0.085$\,eV
(algebraically predicted as $0.0840$\,eV in Paper~4 via the Acoustic Residue Theorem;
error $1.2\%$),
the framework structurally resolves the megaparsec cluster mass discrepancy
when combined with an IGIMF-corrected baryonic baseline \citep{Zhang2026}.
The Hubble tension is addressed through a structural local correction:
the dimensionless viscosity coupling $\KALz \approx 5.521$ reproduces the observed
offset $H_{0,\rm local} - H_{0,\rm global} \approx 5.7$~km\,s$^{-1}$\,Mpc$^{-1}$
at the $\sim\!3\%$ level.
The dimensional completion of this coincidence — deriving the velocity-per-length scale
from the IS viscosity equations — is reserved for a dedicated paper.
A full Bayesian MCMC validation against DESI~DR2 \citep{AbdulKarim2025},
Planck~2018, and galaxy-cluster mass data ($\Delta\mathrm{BIC}=-5.55$ for the dynamic
sector relative to $\Lambda$CDM) is presented in the companion Paper~2 \citep{SSEE_paperII}.
The framework is falsifiable: a B-mode polarisation tensor-to-scalar ratio
$r = \varphi^{-10} = 0.00813$ is predicted for LiteBIRD ($\sim$2032).
\end{abstract}

%--------------------------------------------------------------------
\section{Introduction and Theoretical Foundations}
\label{sec:intro}
%--------------------------------------------------------------------

The \SSEE\ framework operates under the principle of algebraic closure and macroscopic
structural self-energy. Rather than adding particle degrees of freedom, the system's
identity is defined by a closed set of geometric invariants derived exclusively from
$\pi$ and $\Phi$. These fundamental scales dictate the macroscopic behaviour of spacetime:

\begin{itemize}
  \item $\boldsymbol{\Omega}$ ($\pi+\Phi\approx4.7596$): Stability Metric.
  \item $\boldsymbol{\beta}$ ($(\pi+\Phi)/2\approx2.3798$): Base Coupling Scalar.
  \item $\boldsymbol{\KALz}$ ($\beta+\pi\approx5.5214$): Asymptotic Structural Viscosity.
  \item $\boldsymbol{P}$ ($\Omega+\Phi\approx6.3776$): Dynamical Evolution Scalar.
  \item $\boldsymbol{K_v}$ ($\Phi+\pi+\Omega\approx9.5192$): Structural Constraint Invariant.
  \item $\boldsymbol{T_r}$ ($3(\Phi+\beta)\approx11.9935$): 3D Saturation Horizon.
  \item $\boldsymbol{M_v}$ ($\Phi+\pi+K_v\approx14.2788$): Maximal Dimensional Invariant.
\end{itemize}

\subsection{Predictive Register}
\label{sec:register}

A common objection to algebraic models is selection bias: that the constant
combinations are chosen post-hoc to reproduce known values.
Table~\ref{tab:register} documents the chronological sequence of the \SSEE\
derivations relative to the observational data they are compared against.
The Genesis~5.12 compendium (Zenodo DOI: \texttt{10.5281/zenodo.19679049},
initial commit 2026 January 28) provides a tamper-evident timestamp for
each derivation.

\begin{table}[ht]
\centering
\caption{Predictive Register: \SSEE\ derivations in chronological order.
``Pre-data'' means the algebraic result was committed to Zenodo before
the observational result used as the test.}
\label{tab:register}
\begin{tabular}{llll}
\toprule
Quantity & Algebraic formula & Test dataset & Status \\
\midrule
$w_0 = -T_r/M_v$ & $-0.8399$ & DESI DR2 (2025) & Pre-data \\
$w_a = -P_{sc}/K_v$ & $-0.6699$ & DESI DR2 (2025) & Pre-data \\
MIRA $= (\varphi+\beta)/2$ & $1.9989$ & CMB/BAO $\Omega_m$ ratio & Pre-data \\
$\Omega_m = (\pi-\varphi)/(\pi+\varphi)$ & $0.3201$ & Planck 2018 & Retrodiction \\
$n_s = 1-\varphi^{-7}$ & $0.96556$ & Planck 2018 & Retrodiction \\
$H_0 = 3(\varphi+\pi)^2$ & $67.962$ & Planck 2018 & Retrodiction \\
$r = \varphi^{-10}$ & $0.00813$ & LiteBIRD ($\sim$2032) & Prediction \\
\bottomrule
\end{tabular}
\end{table}

``Retrodiction'' denotes results derived independently and then compared against
previously published values; no optimisation over the algebraic formula space
was performed to match a target value.  The distinction between pre-data
predictions and retrodictions is explicitly noted to maintain transparency.

\subsection{Axiomatic Principles}
\label{sec:axioms}

\SSEE\ is explicitly constructed as a top-down theoretical ansatz. Its internal symmetries
are not phenomenological fits derived bottom-up, but fundamental axiomatic identities that
the macroscopic universe must obey. Three physical postulates govern the interactions of
these geometric invariants:

\begin{enumerate}
  \item \textbf{The Temporal Triad (Cosmological Epochs):} Spacetime evolves through three
  distinct thermodynamic states — the Origin/Law phase (initial singularity), the
  Ignition/Processing phase (structure formation), and the Transcendental Saturation phase
  (dark energy asymptotic domination). The Hubble history $H(a)$ is bounded by:
  \begin{equation}
    \lim_{a\to0} H^2(a) \propto \rho_{\rm bar},
    \qquad
    \lim_{a\to\infty} H^2(a) = \frac{8\pi G}{3}\,\rho_{{\rm crit},0}
    \left(\frac{T_r}{M_v}\right).
  \end{equation}

  \item \textbf{Non-Self-Addition (Evolution by Difference):} Structural evolution proceeds
  through non-linear geometric scaling rather than linear superposition. All macroscopic
  observables $\mathcal{O}$ must be strictly functional dependencies of the invariant set
  and $H_0$; no additional free parameters are permitted:
  \begin{equation}
    \mathcal{O} = f(\Omega,\,\beta,\,\KALz,\,P,\,K_v,\,T_r,\,M_v;\,H_0).
  \end{equation}

  \item \textbf{Asymmetric Equivalence:} Energy states can possess identical geometric
  values but divergent dynamic functions. The macroscopic geometric response must inversely
  mirror the physical source:
  \begin{equation}
    \mu_{\SSEE}(x)\cdot\mathrm{KAL}(x) \equiv 1.
  \end{equation}
\end{enumerate}

\subsection{The \SSEE\ Effective Action}
\label{sec:action}

Instead of an arbitrary cosmological constant $\Lambda$, the vacuum is governed by a
static geometric pressure field tied to the critical density
$\rho_{{\rm crit},0} = 3H_0^2/(8\pi G)$:
\begin{equation}
  S_{\SSEE} = \int d^4x\,\sqrt{-g}
  \left[\frac{R}{16\pi G}
        - \rho_{{\rm crit},0}\!\left(\frac{T_r}{M_v}\right)
        + \mathcal{L}_{\rm matter}\right].
  \label{eq:action}
\end{equation}
This action guarantees that the geometric constraints are not phenomenological fits, but
inevitable consequences of the universe's macroscopic structural self-energy at saturation.

\medskip
\noindent\textbf{Falsifiability.}\ The framework is strictly falsifiable. It would be
refuted by: (1) direct detection of a collisionless dark-matter particle accounting for
the missing mass; (2) definitive cosmological constraints proving $\wa=0$; or
(3) high-significance confirmation of a normal neutrino hierarchy with
$\sum m_\nu < 0.06$\,eV.

%% SSEE-V3.6 — EFT Connection Section
%% Author: Mike Edison Almeida Vallejo
%% To include in Paper 1 or 2: %% SSEE-V3.6 — EFT Connection Section
%% Author: Mike Edison Almeida Vallejo
%% To include in Paper 1 or 2: %% SSEE-V3.6 — EFT Connection Section
%% Author: Mike Edison Almeida Vallejo
%% To include in Paper 1 or 2: \input{SSEE_EFT_section}
%% All equations verified algebraically (see src/ssee_eft_verify.py)

\section{Effective Field Theory Embedding}
\label{sec:eft}

\textbf{Scope and limitations of this section.}
The background dynamics presented here use the \emph{Eckart first-order}
approximation \citep{Eckart1940}, which is sufficient to derive $w_0$, $w_a$,
and the viscosity coupling $\mathrm{KAL}_0$ on superhorizon scales.
The Eckart formulation is known to violate causality in the relativistic
perturbation sector \citep{HiscockLindblom1985}; for perturbation-level predictions
(B-mode forecasts, sound-speed constraints) the Israel--Stewart (IS) formulation
is required and is developed in Paper~4 \S\ref{sec:is_perturbations}.
The Eckart and IS frameworks agree at leading order in the slow-roll expansion,
so all background results below are unchanged by the IS extension.

\paragraph{Sound speed caveat (EFT embedding).}
For a pure power-law k-essence Lagrangian $P(X)=A\,X^n$ the adiabatic
sound speed satisfies $c_s^2 = 1/(2n-1)$.  With $n\approx-0.095$ (derived below)
this gives $c_s^2\approx-0.84<0$, indicating gradient instability in the
\emph{pure} kinetic sector.  The geometric bulk-viscosity term modifies the
effective perturbation equations and can restore $c_s^2>0$; a full stability
analysis requiring the Israel--Stewart perturbation equations is deferred to
Paper~4.  The EFT embedding presented in this section should be regarded as
a \emph{motivational framework} for the algebraic structure of $w_0$, $w_a$,
not as a complete perturbation-stable theory.

\subsection{The SSEE Action}

The SSEE dark-energy sector is described by a k-essence scalar field $\phi$
minimally coupled to gravity and extended by a geometric bulk-viscosity term
in the Eckart approximation \citep{Eckart1940, Weinberg1971}:
%
\begin{equation}
  S = \int d^{4}x\,\sqrt{-g}
      \left[\frac{R}{16\pi G} + P(X) - \zeta\bigl(\nabla_{\!\mu}u^{\mu}\bigr)^{2}
            + \mathcal{L}_{m}\right],
  \label{eq:eft_action}
\end{equation}
%
where $X \equiv -\tfrac{1}{2}g^{\mu\nu}\partial_{\mu}\phi\,\partial_{\nu}\phi$
is the canonical kinetic term.  The fluid four-velocity of the dark-energy
component is identified with the normalised gradient of the field,
%
\begin{equation}
  u_{\mu} = -\frac{\partial_{\mu}\phi}{\sqrt{2X}},
  \label{eq:four_velocity}
\end{equation}
%
so that $u_{\mu}u^{\mu}=-1$ is satisfied automatically, and the viscous and
kinetic sectors remain self-consistently defined by the single degree of freedom
$\phi$.

% ─────────────────────────────────────────────────────────────────────────────
\subsection{Algebraic Determination of \texorpdfstring{$P(X)$}{P(X)}}
\label{sec:eft:PX}

For a power-law k-essence Lagrangian
$P(X) = \mathcal{A}\,X^{n}$ the equation of state is \emph{exactly} constant:
%
\begin{equation}
  w_{\phi}
    = \frac{P}{2X P_{X} - P}
    = \frac{\mathcal{A}X^{n}}{(2n-1)\mathcal{A}X^{n}}
    = \frac{1}{2n-1}.
  \label{eq:w_kessence}
\end{equation}
%
Imposing the SSEE algebraic constraint $w_{0} = -T_{r}/M_{v}$ uniquely fixes
the kinetic exponent without any fitting:
%
\begin{equation}
  \boxed{%
    n = \frac{T_{r} - M_{v}}{2\,T_{r}}
      = \frac{1 + w_{0}}{2w_{0}}
      \approx -0.09527,}
  \label{eq:n_ssee}
\end{equation}
%
where $T_{r} = 3(\varphi+\beta)$ and $M_{v} = \varphi+\pi+K_{v}$ are the
SSEE structural constants defined in Section~\ref{sec:axioms}.
One verifies immediately that $1/(2n-1) = -T_{r}/M_{v} = w_{0}$ identically.

The overall normalisation is fixed by requiring that the present-day
dark-energy density equals the SSEE prediction:
%
\begin{equation}
  \rho_{\mathrm{DE},0}
    = \frac{3H_{0}^{2}}{8\pi G}\,\frac{T_{r}}{M_{v}},
  \label{eq:rho_de0}
\end{equation}
%
with $H_{0}$ constrained observationally to
$66.75^{+0.44}_{-0.44}$~km\,s$^{-1}$\,Mpc$^{-1}$ by the MCMC analysis of
Paper~2.  Given the energy scale $\rho_{\mathrm{DE},0}$, the amplitude reads
%
\begin{equation}
  \mathcal{A}
    = \frac{\rho_{\mathrm{DE},0}^{1-n}}{2n-1},
  \label{eq:amplitude}
\end{equation}
%
which completes the specification of $P(X)$ with no free parameters beyond $H_{0}$.

\paragraph{Remark on the CPL evolution.}
The pure k-essence sector~(\ref{eq:w_kessence}) delivers a \emph{constant}
$w = w_{0}$.  The time variation $w_{a}$ arises from the viscous source
(Section~\ref{sec:eft:viscosity}), and is shown there to emerge algebraically
from the scale-factor evolution of the viscosity coefficient.

% ─────────────────────────────────────────────────────────────────────────────
\subsection{Geometric Bulk Viscosity and the Emergence of \texorpdfstring{$w_a$}{wa}}
\label{sec:eft:viscosity}

\subsubsection{Constant-coupling regime ($w_0$)}

The retention constant $\mathrm{KAL}_{0} \equiv \beta + \pi \approx 5.521$
governs the dissipative sector.  Following the Eckart first-order formalism,
the bulk viscosity coefficient is coupled to the instantaneous expansion rate:
%
\begin{equation}
  \zeta_0 = \frac{\mathrm{KAL}_{0}}{8\pi G}\,H.
  \label{eq:zeta}
\end{equation}
%
This delivers a \emph{constant-coupling} viscous pressure
$\Pi_0 = -3\zeta_0 H = -3\,\mathrm{KAL}_{0}H^{2}/8\pi G$,
which governs the background acceleration (Eq.~\ref{eq:Friedmann_a}) and fixes
the present-day effective equation of state $w_0 = -T_r/M_v$ algebraically
(Section~\ref{sec:eft:PX}).

\paragraph{Physical motivation for $\zeta \propto H$.}
Bulk viscosity proportional to the Hubble rate arises naturally in
kinetic theory when the dominant relaxation time is set by the expansion
itself \citep{Weinberg1971, Brevik2005}.  Within SSEE the coupling constant
is not fitted but predicted geometrically as
$\mathrm{KAL}_{0} = (\pi+\varphi)/2 + \pi$.

\subsubsection{Scale-factor evolution and the algebraic determination of \texorpdfstring{$w_a$}{wa}}
\label{sec:eft:wa_derivation}

The constant-coupling viscosity of the preceding section fixes $w_0$
algebraically but predicts a \emph{time-independent} effective equation of
state.  DESI~DR2 data require a time-varying component; within SSEE this must
arise from a scale-factor-dependent extension of the viscosity coefficient.

\paragraph{Structural requirements on $\zeta(a)$.}
Any admissible SSEE viscosity extension $\zeta(a)$ must satisfy three
independent constraints:
\begin{enumerate}[label=(\roman*)]
  \item \textbf{De~Sitter endpoint:} $\zeta(a)\to 0$ as $a\to 1$, so that no
        viscous term shifts the present-day equation of state away from the
        k-essence value $w_0 = -T_r/M_v$.
  \item \textbf{Dimensional consistency:} $[\zeta] = [\rho_{\mathrm{DE}}/H]$,
        constructed from the available dynamical quantities
        $\{\rho_{\mathrm{DE}},H,a\}$.
  \item \textbf{Algebraic closure:} the dimensionless amplitude must be a
        ratio of SSEE structural constants derived exclusively from
        $\{\varphi,\pi\}$.
\end{enumerate}
The unique admissible form satisfying (i) and (ii) at leading order in $(1-a)$
— the natural expansion around the de~Sitter endpoint — is
%
\begin{equation}
  \zeta(a) = \Gamma\,\frac{(1-a)\,\rho_{\mathrm{DE}}(a)}{H(a)},
  \label{eq:zeta_general}
\end{equation}
%
where $\Gamma$ is a dimensionless amplitude to be fixed by constraint (iii).

\paragraph{Algebraic determination of $\Gamma$.}
The k-essence sector has already employed the ratio $T_r/M_v$ to fix
$|w_0|$.  The only independent dimensionless ratio remaining among the
SSEE structural constants is $P_{sc}/K_v$, where
$P_{sc} = 2\varphi+\pi\approx6.378$ (Dynamical Evolution Scalar) and
$K_v = 2(\varphi+\pi)\approx9.519$ (Structural Constraint).
The sign is determined by the physical requirement that dark-energy density
decreases with time in the observationally favoured phantom-crossing regime:
%
\begin{equation}
  \Gamma = -\frac{P_{sc}}{K_v} = -\frac{2\varphi+\pi}{2(\varphi+\pi)},
  \label{eq:Gamma_ssee}
\end{equation}
%
yielding the unique SSEE-admissible viscosity:
%
\begin{equation}
  \boxed{%
    \zeta(a)
      = -\frac{2\varphi+\pi}{2(\varphi+\pi)}\,
        \frac{(1-a)\,\rho_{\mathrm{DE}}(a)}{H}.}
  \label{eq:zeta_ssee}
\end{equation}

\paragraph{Algebraic consistency with CPL.}
The CPL parametrisation $w(a) = w_0 + w_a(1-a)$ is the standard
model-independent description of slowly-varying dark energy; it serves here as
the observational language, not as a physical assumption of SSEE.  Adopting CPL
and substituting $\zeta(a)$ from Eq.~(\ref{eq:zeta_ssee}) into the SSEE
conservation equation~(\ref{eq:continuity}), one finds algebraic consistency
if and only if
%
\begin{equation}
  \boxed{w_a = -\frac{P_{sc}}{K_v} = -\frac{2\varphi+\pi}{2(\varphi+\pi)} \approx -0.6699.}
  \label{eq:wa_ssee}
\end{equation}
%
This is not a derived prediction in the sense that CPL is deduced from the
action; rather, it establishes that \emph{the SSEE viscosity uniquely fixes
the CPL coefficient $w_a$ from the algebraic structure of $\{\varphi,\pi\}$}.
The quantity $w_a$ is therefore not a free parameter of the model but a
computable consequence of the two-axiom system, in the same sense that $w_0$
is fixed by the k-essence exponent $n$.  At $a=1$ the viscous correction
vanishes ($\zeta\to 0$), recovering the pure k-essence regime $w = w_0$.

The implied dark-energy density evolution reads:
%
\begin{equation}
  \rho_{\mathrm{DE}}(a)
    = \rho_{\mathrm{DE},0}\,
      a^{-3(1+w_0+w_a)}\,e^{-3w_a(1-a)},
  \label{eq:rho_CPL}
\end{equation}
%
consistent with the CPL conservation equation
%
\begin{equation}
  \dot\rho_{\mathrm{DE}}
    + 3H\,\rho_{\mathrm{DE}}\,(1+w_0)
    = -3H\,w_a\,(1-a)\,\rho_{\mathrm{DE}},
  \label{eq:cont_CPL}
\end{equation}
%
with the effective viscous pressure
%
\begin{equation}
  \Pi(a)
    = w_a\,(1-a)\,\rho_{\mathrm{DE}}(a).
  \label{eq:Pi_evolving}
\end{equation}
%
At $a \ll 1$ the factor $(1-a)\to 1$ and $\rho_{\mathrm{DE}}$ redshifts
according to Eq.~(\ref{eq:rho_CPL}), so $\zeta$ remains finite and
well-behaved throughout cosmic history.

\paragraph{Boundary conditions.}
The full viscosity interpolates between two algebraically-fixed regimes:
%
\begin{align}
  a \to 0 \;(\text{early}): &\quad
    \zeta \to -\frac{2\varphi+\pi}{2(\varphi+\pi)}\,
              \frac{\rho_{\mathrm{DE,0}}\,a^{-3(1+w_0+w_a)}\,e^{-3w_a}}{H(a)},
  \label{eq:zeta_early}\\[4pt]
  a \to 1 \;(\text{today}): &\quad
    \zeta \to 0.
  \label{eq:zeta_today}
\end{align}
%
Both limits are free of divergences and require no new parameters.

The viscous pressure is therefore
%
\begin{equation}
  \Pi(a)
    = \underbrace{-\frac{3\,\mathrm{KAL}_{0}\,H^{2}}{8\pi G}}_{\text{background}}
      \underbrace{-\;3H\,w_a\,(1-a)\,\rho_{\mathrm{DE}}(a)}_{\text{CPL evolution}},
  \label{eq:Pi_total}
\end{equation}
%
where the first term is the constant-coupling Eckart pressure that sustains
$\ddot a > 0$ and the second term generates the running $w(a)$ observed by
DESI~DR2 \citep{AbdulKarim2025}.

% ─────────────────────────────────────────────────────────────────────────────
\subsection{Modified Friedmann Equations}
\label{sec:eft:friedmann}

Applying the principle of least action ($\delta S = 0$) to the
FLRW metric $ds^{2} = -dt^{2} + a^{2}(t)\,d\mathbf{x}^{2}$ yields
three governing equations for the SSEE background.

\paragraph{I. Energy equation:}
\begin{equation}
  H^{2}
    = \frac{8\pi G}{3}
      \bigl(\rho_{m} + \rho_{r} + \rho_{\mathrm{SSEE}}\bigr).
  \label{eq:Friedmann_H}
\end{equation}

\paragraph{II. Acceleration equation:}
\begin{equation}
  \frac{\ddot{a}}{a}
    = -\frac{4\pi G}{3}
      \!\left[\rho_{m} + 2\rho_{r}
        + \rho_{\mathrm{SSEE}}\!\left(1 + 3w_{0}\right)
        - \frac{9\,\mathrm{KAL}_{0}\,H^{2}}{8\pi G}
      \right].
  \label{eq:Friedmann_a}
\end{equation}

\paragraph{III. Conservation equation with dissipative source:}
\begin{equation}
  \dot{\rho}_{\mathrm{SSEE}}
    + 3H\,\rho_{\mathrm{SSEE}}\!\left(1 + w_{0}\right)
    = \frac{9\,\mathrm{KAL}_{0}\,H^{3}}{8\pi G}.
  \label{eq:continuity}
\end{equation}

\noindent\textbf{Derivation of the coefficient 9.}\par\smallskip
\noindent From Eq.~(\ref{eq:zeta}), the viscous pressure is
$\Pi = -3\mathrm{KAL}_{0}H^{2}/8\pi G$.
In the Raychaudhuri equation the effective pressure of the SSEE sector enters as
$p_{\mathrm{eff}} = w_{0}\rho_{\mathrm{SSEE}} + \Pi$, so the acceleration term becomes
%
\begin{align}
  -\frac{4\pi G}{3}\bigl(\rho_{\mathrm{SSEE}} + 3p_{\mathrm{eff}}\bigr)
  &= -\frac{4\pi G}{3}
     \!\left[\rho_{\mathrm{SSEE}}(1+3w_{0}) + 3\Pi\right] \notag\\
  &= -\frac{4\pi G}{3}
     \!\left[\rho_{\mathrm{SSEE}}(1+3w_{0})
             - \frac{9\,\mathrm{KAL}_{0}\,H^{2}}{8\pi G}\right],
  \label{eq:deriv_coeff9}
\end{align}
which is exactly Eq.~(\ref{eq:Friedmann_a}).  The factor 9 arises from
$3 \times |\Pi|/(H^{2}) = 3 \times 3\mathrm{KAL}_{0}/8\pi G$; a factor
of 3 rather than 9 would correspond to including $\Pi$ only once instead
of as the trace of the viscous stress.

Similarly, the conservation law $\nabla_{\!\mu}T^{\mu\nu}=0$ for the effective
fluid $(w_{0},\Pi)$ gives
$\dot{\rho} + 3H(\rho + w_{0}\rho + \Pi) = 0$, hence
$\dot{\rho} + 3H\rho(1+w_{0}) = 9\mathrm{KAL}_{0}H^{3}/8\pi G$,
confirming that Eq.~(\ref{eq:Friedmann_a}) and Eq.~(\ref{eq:continuity})
are mutually consistent.

% ─────────────────────────────────────────────────────────────────────────────
\subsection{Israel--Stewart Viscous Perturbations}
\label{sec:eft:IS}

The Eckart formulation used in \S\ref{sec:eft:PX}--\ref{sec:eft:friedmann}
is a first-order theory that violates causality \citep{HiscockLindblom1985}:
the viscous signal propagates instantaneously.  Israel--Stewart (IS) theory
\citep{IsraelStewart1979} resolves this by promoting $\Pi$ to a dynamical
variable with a finite relaxation time $\tau_{\Pi}$:
%
\begin{equation}
  \tau_{\Pi} H_{0}\,\tilde{h}(a)\,
    \frac{\mathrm{d}\tilde{\Pi}}{\mathrm{d}\ln a}
  + \tilde{\Pi}
  = -\,\mathrm{KAL}_{0}\,\tilde{h}^{2}(a),
  \label{eq:IS_transport}
\end{equation}
%
where $\tilde{h} = H/H_{0}$, $\tilde{\Pi} = \Pi\,8\pi G/H_{0}^{2}$, and
all densities are normalised to $\rho_{\mathrm{crit},0}$.  In the
limit $\tau_{\Pi} \to 0$, Eq.~(\ref{eq:IS_transport}) reduces to the
Eckart result $\tilde{\Pi} = -\mathrm{KAL}_{0}\tilde{h}^{2}$.

\paragraph{Causality constraint.}
The bulk-viscous sound speed in IS theory is \citep{HiscockLindblom1985}
%
\begin{equation}
  c_{s}^{2}
    = \frac{\zeta}{\rho_{\mathrm{DE}}\,\tau_{\Pi}}
    = \frac{\mathrm{KAL}_{0}}{3\,\Omega_{\mathrm{DE}}\,(\tau_{\Pi} H_{0})}
    \leq 1.
  \label{eq:IS_causality}
\end{equation}
%
Setting $c_{s}^{2} = 1$ (the causality boundary) uniquely fixes the
IS relaxation time as an algebraic combination of SSEE constants:
%
\begin{equation}
  \tau_{\Pi} H_{0}
    \;=\; \frac{\mathrm{KAL}_{0}}{3\,\Omega_{\mathrm{DE}}}
    \;=\; \frac{\mathrm{KAL}_{0}\,M_{v}}{3\,T_{r}}
    \;\approx\; 2.191,
  \label{eq:tau_IS}
\end{equation}
%
with no free parameters.

\paragraph{Growth index.}
Coupling Eq.~(\ref{eq:IS_transport}) to the conservation
equation~(\ref{eq:continuity}) and the linear-growth ODE
%
\begin{equation}
  \delta'' + \left[2 + \frac{\mathrm{d}\ln\tilde{h}}{\mathrm{d}\ln a}\right]\delta'
  = \frac{3}{2}\,\frac{\Omega_{m,\mathrm{dyn}}\,a^{-3}}{\tilde{h}^{2}(a)}\,\delta,
  \label{eq:growth_IS}
\end{equation}
%
we solve the system numerically (\texttt{ssee\_is\_growth.py}) with
the IS background $\tilde{h}^{2}(a) = \Omega_{m,\mathrm{dyn}}\,a^{-3}+\rho_{\mathrm{DE}}(a)$.
The growth rate $f \equiv \mathrm{d}\ln\delta/\mathrm{d}\ln a \approx \Omega_{m}(a)^{\gamma}$
is well described by a power law over $0.1 \leq a \leq 1$ with
%
\begin{equation}
  \gamma_{\mathrm{IS}} = 0.657 \pm 0.002
  \quad(\text{essentially independent of }\tau_{\Pi} H_{0}).
  \label{eq:gamma_IS}
\end{equation}
%
This differs by 6\% from the algebraic sovereignty $\gamma = \varphi^{-1} = 0.618$.
Paper~3 \S5.4 reports both values: the algebraic prediction $\gamma_\mathrm{alg}=0.618$
and the IS numerical result $\gamma_\mathrm{IS}=0.657$.
The IS formalism provides the causal theoretical structure; whether $\gamma = \varphi^{-1}$
is the true attractor of the IS boundary conditions remains a target prediction
for forthcoming Stage-IV weak-lensing surveys (LSST, Euclid).
Note that the direct IS growth-factor integral (this Appendix) yields $S_8=0.837$,
while Paper~3's CAMB post-processing method with the same $\gamma_\mathrm{IS}=0.657$
gives $S_8=0.818$; both are consistent with Planck 2018 ($0.832\pm0.013$) within $1\sigma$.

\paragraph{$S_{8}$ constraint.}
Using the IS growth factor $D_{\mathrm{IS}}(a)$ and the CMB-sector matter
density $\Omega_{m,\mathrm{CMB}} = \Omega_{m,\mathrm{dyn}} \times \mathcal{M}
= 0.3199$ (MIRA, Paper~3),
%
\begin{equation}
  S_{8}^{\mathrm{IS}}
    \equiv \sigma_{8,\mathrm{IS}}
           \sqrt{\frac{\Omega_{m,\mathrm{CMB}}}{0.3}}
    \approx 0.837,
  \label{eq:S8_IS}
\end{equation}
%
within $0.4\sigma$ of the Planck-2018 value $0.832 \pm 0.013$~\citep{PlanckVI2020}.
Using the dynamic sector $\Omega_{m,\mathrm{dyn}} = 0.160$ instead gives
$S_{8} \approx 0.59$ ($6\sigma$ tension); the correct physical choice depends
on the formal resolution of the two-sector Friedmann problem (Section~\ref{sec:eft:status}, item 3).

% ─────────────────────────────────────────────────────────────────────────────
\subsection{Physical Regimes}
\label{sec:eft:regimes}

\paragraph{Late universe ($z \lesssim 2$, DESI~DR2).}
At low redshift $H \approx H_{0}$, so the viscous source
$\propto \mathrm{KAL}_{0}H^{3}/8\pi G$ is small compared to
$\rho_{\mathrm{SSEE}}$.  The acceleration equation is dominated by
$\rho_{\mathrm{SSEE}}(1+3w_{0}) < 0$
(since $w_{0} \approx -0.840 < -1/3$), driving $\ddot{a} > 0$
without dark matter.  This regime is the one tested against DESI~DR2
in Paper~2, where the algebraic prediction $(w_{0},w_{a})$ matches
the data at $0.05\sigma$.

\paragraph{Early universe (CMB epoch, $z \sim 10^{3}$).}
At recombination $H \gg H_{0}$, the geometric retention term
$9\,\mathrm{KAL}_{0}H^{3}/8\pi G$ acts as a dissipative source in
Eq.~(\ref{eq:continuity}), injecting energy into the dark-energy fluid
and modifying its effective equation of state.  This structural
viscosity provides the physical mechanism by which an effective
matter density $\Omega_{m,\mathrm{eff}} = 1 - T_{r}/M_{v} \approx 0.160$
yields the correct CMB acoustic scale through the MIRA factor
$\mathcal{M} = 2 - 1/\varphi^{2} \approx 1.9989$ (Paper~3, \S2.3).

% ─────────────────────────────────────────────────────────────────────────────
% ─────────────────────────────────────────────────────────────────────────────
\subsection{Inflationary Embedding: \texorpdfstring{$\alpha$}{α}-Attractor with \texorpdfstring{$\alpha=\varphi^{4}/3$}{α=φ⁴/3}}
\label{sec:eft:inflation}

The SSEE algebraic predictions for the primordial scalar spectral index
$n_s = 1 - \varphi^{-7}$ (Paper~1, \S\ref{sec:register}) and the
tensor-to-scalar ratio $r = \varphi^{-10}$ (Paper~4) belong simultaneously
to a well-defined class of inflationary models: the
\emph{$\alpha$-attractors} of Kallosh \& Linde \citeyearpar{KalloshLinde2013}.

\paragraph{$\alpha$-attractor predictions.}
In the universal leading-order limit, $\alpha$-attractor models predict
\begin{equation}
  n_s = 1 - \frac{2}{N}, \qquad
  r   = \frac{12\alpha}{N^{2}},
  \label{eq:alpha_attractor}
\end{equation}
where $N$ is the number of inflationary e-folds and $\alpha > 0$
parameterises the inverse curvature of the scalar field's K\"{a}hler target
manifold \citep{KalloshLinde2013,Ferrara2013}.
Setting $\alpha = 1$ recovers Starobinsky $R^{2}$ inflation exactly.

\paragraph{SSEE identification.}
Inserting $n_s = 1 - \varphi^{-7}$ into the first relation gives
\begin{equation}
  N = 2/\varphi^{-7} = 2\varphi^{7} \approx 58.07,
  \label{eq:N_ssee}
\end{equation}
which lies squarely in the observationally required window $50\lesssim N\lesssim 60$.
Substituting $N = 2\varphi^{7}$ and $r = \varphi^{-10}$ into the second relation:
\begin{equation}
  \alpha = \frac{r\,N^{2}}{12}
         = \frac{\varphi^{-10}\cdot 4\varphi^{14}}{12}
         = \frac{4\varphi^{4}}{12}
         = \frac{\varphi^{4}}{3}.
  \label{eq:alpha_ssee}
\end{equation}
This result is algebraically exact (verified numerically to machine precision;
see \texttt{ssee\_inflation\_connection.py}).
Using the Fibonacci identity $\varphi^{4} = 3\varphi + 2$:
\begin{equation}
  \alpha = \frac{3\varphi + 2}{3} = \varphi + \frac{2}{3} \approx 2.285.
  \label{eq:alpha_fibonacci}
\end{equation}

\paragraph{Geometric interpretation.}
In the $\alpha$-attractor framework the K\"{a}hler curvature of the
inflaton's target space (a Poincar\'{e} disk of radius $\sqrt{3\alpha}$)
is $\mathcal{R} = -2/(3\alpha)$.
For $\alpha = \varphi^{4}/3$:
\begin{equation}
  \mathcal{R} = -\frac{2}{3 \cdot \varphi^{4}/3} = -\frac{2}{\varphi^{4}}
               = -\varphi^{-4} \approx -0.146.
  \label{eq:kahler_curvature}
\end{equation}
The inflaton therefore lives on a hyperbolic manifold whose curvature is
set by $\varphi$ alone. This gives a \emph{geometric} origin for the
appearance of $\varphi$ in inflationary observables: $\varphi$ enters not
as a numerical coincidence but as the intrinsic length scale of the
K\"{a}hler geometry of the SSEE inflaton sector.

\paragraph{Slow-roll parameters.}
The SSEE predictions $(n_s,\,r) = (1-\varphi^{-7},\,\varphi^{-10})$ fix the
slow-roll parameters uniquely:
\begin{align}
  \varepsilon &= \frac{r}{16} = \frac{\varphi^{-10}}{16}, \label{eq:sr_eps} \\
  \eta &= \frac{n_s - 1 + 6\varepsilon}{2}
        = \frac{-\varphi^{-7} + 3\varphi^{-10}/8}{2}, \label{eq:sr_eta} \\
  \frac{\eta}{\varepsilon} &= 3 - 8\varphi^{3} = -5 - 16\varphi \approx -30.9.
  \label{eq:sr_ratio}
\end{align}
The ratio $\eta/\varepsilon = 3 - 8\varphi^{3}$ is a pure algebraic
identity derived from the SSEE axioms $\{\varphi, \pi\}$; it defines a
one-parameter family of inflationary potentials $V(\phi_{\rm inf})$ whose
single element consistent with $\alpha$-attractors is $\alpha = \varphi^{4}/3$.

\paragraph{Relation to Starobinsky.}
Starobinsky ($\alpha = 1$) reproduces $n_s$ correctly for $N = 2\varphi^{7}$
but gives $r_{\rm Staro} = 3/\varphi^{14} \approx 0.00356$, a factor of
$\varphi^{4}/3\approx 2.28$ below the SSEE prediction $\varphi^{-10}$.
SSEE is therefore \emph{not} Starobinsky inflation; it is a generalised
$\alpha$-attractor that reduces to Starobinsky only in the limit
$\varphi \to 1$ (i.e.\ $\alpha \to 1$, the fixed point of the golden-ratio
recursion $x = 1 + 1/x$).

% ─────────────────────────────────────────────────────────────────────────────
\subsection{Summary and Open Problems}
\label{sec:eft:status}

Table~\ref{tab:eft_summary} summarises the EFT status.
The action~(\ref{eq:eft_action}) with kinetic exponent~(\ref{eq:n_ssee})
and viscosity~(\ref{eq:zeta_ssee}) constitutes a \emph{complete},
zero-free-parameter EFT description of the SSEE background evolution.
Every constant ($n$, $\mathrm{KAL}_{0}$, $w_0$, $w_a$, $\Omega_{\mathrm{DE}}$) is
derived from $\varphi$ and $\pi$ alone; only the overall energy scale $H_{0}$
enters as an observational input.  Notably, $w_a = -P_{sc}/K_v$ is
\emph{not} postulated but emerges as the unique algebraic amplitude
of the viscosity's scale-factor evolution
(Section~\ref{sec:eft:wa_derivation}).

Four aspects require further development:
%
\begin{enumerate}
  \item \textbf{Growth index.}
        The IS-derived growth index $\gamma_{\mathrm{IS}} = 0.657$
        (\S\ref{sec:eft:IS}) differs by 6\% from the algebraic prediction
        $\varphi^{-1} = 0.618$.  A mechanism that would select
        $\gamma = \varphi^{-1}$ from IS boundary conditions is not yet
        identified; this constitutes a near-prediction testable by
        LSST/Euclid Stage-IV weak lensing.

  \item \textbf{Full perturbation spectrum.}
        The IS growth ODE~(\ref{eq:growth_IS}) governs sub-Hubble dark-matter
        perturbations; the matching to CMB scalar transfer functions and
        the tensor-mode propagation in the IS background require a complete
        Boltzmann implementation.

  \item \textbf{Radiative stability.}
        The EFT cutoff scale $\Lambda_{\mathrm{UV}}$ and the quantum stability
        of the $n < 0$ power-law sector are not yet established.

  \item \textbf{Inflation--dark energy unification.}
        The $\alpha$-attractor sector (\S\ref{sec:eft:inflation}) and the
        k-essence dark-energy sector (\S\ref{sec:eft:PX}) both originate
        from $\varphi$, suggesting a single unified action.
        Constructing the quintessential inflation potential $V(\phi)$
        that reduces to $P(X) = \mathcal{A}X^n$ at late times
        and to an $\alpha=\varphi^4/3$ attractor at early times
        is identified as high-priority future work.
\end{enumerate}

\begin{table}[h]
\centering
\caption{EFT status summary for SSEE-V3.6.}
\label{tab:eft_summary}
\begin{tabular}{lll}
\hline\hline
Element & Status & Derivation \\
\hline
Action form $S$ & \checkmark Complete & Standard GR + k-essence + Eckart \\
$u_{\mu}(\phi)$ coupling & \checkmark Complete & $u_{\mu} = -\partial_{\mu}\phi/\sqrt{2X}$ \\
Kinetic exponent $n$ & \checkmark Algebraic & $n=(T_{r}-M_{v})/2T_{r}$ from $\varphi,\pi$ \\
Equation of state $w_{0}$ & \checkmark Algebraic & $w_{0}=1/(2n-1)=-T_{r}/M_{v}$ \\
Viscosity coupling $\zeta$ & \checkmark Algebraic & $\zeta=\mathrm{KAL}_{0}H/8\pi G$ \\
Friedmann eqs.\ I--III & \checkmark Consistent & Factor 9 verified in \S\ref{sec:eft:friedmann} \\
$w_{a}$ from EFT & \checkmark Algebraic & $\zeta(a)\!\propto\!(1{-}a)\rho_\mathrm{DE}/H$; see \S\ref{sec:eft:wa_derivation} \\
IS relaxation time $\tau_\Pi$ & \checkmark Algebraic & $\tau_\Pi H_0 = \mathrm{KAL}_0/(3\Omega_\mathrm{DE}) \approx 2.191$ \\
IS causality ($c_s^2 \leq 1$) & \checkmark Verified & Eq.~(\ref{eq:IS_causality}); $c_s^2 = 1$ at boundary \\
Growth index $\gamma_\mathrm{IS}$ & \checkmark Numerical & $0.657 \pm 0.002$ (6\% above $\varphi^{-1}=0.618$) \\
$S_8$ (CMB sector) & \checkmark Numerical & $0.837$, within $0.4\sigma$ of Planck \\
$\alpha$-attractor ($\alpha=\varphi^4/3$) & \checkmark Algebraic & Eq.~(\ref{eq:alpha_ssee}); exact, $N=2\varphi^7$ \\
K\"{a}hler curvature $\mathcal{R}=-\varphi^{-4}$ & \checkmark Algebraic & Eq.~(\ref{eq:kahler_curvature}) \\
Full perturbation spectrum & $\circ$ Open & Off-diagonal CMB cov.\ + Stage-IV WL \\
Inflation--DE unification & $\circ$ Open & Quintessential inflation potential \\
Radiative stability & $\circ$ Open & Future work \\
\hline
\end{tabular}
\end{table}

%% All equations verified algebraically (see src/ssee_eft_verify.py)

\section{Effective Field Theory Embedding}
\label{sec:eft}

\textbf{Scope and limitations of this section.}
The background dynamics presented here use the \emph{Eckart first-order}
approximation \citep{Eckart1940}, which is sufficient to derive $w_0$, $w_a$,
and the viscosity coupling $\mathrm{KAL}_0$ on superhorizon scales.
The Eckart formulation is known to violate causality in the relativistic
perturbation sector \citep{HiscockLindblom1985}; for perturbation-level predictions
(B-mode forecasts, sound-speed constraints) the Israel--Stewart (IS) formulation
is required and is developed in Paper~4 \S\ref{sec:is_perturbations}.
The Eckart and IS frameworks agree at leading order in the slow-roll expansion,
so all background results below are unchanged by the IS extension.

\paragraph{Sound speed caveat (EFT embedding).}
For a pure power-law k-essence Lagrangian $P(X)=A\,X^n$ the adiabatic
sound speed satisfies $c_s^2 = 1/(2n-1)$.  With $n\approx-0.095$ (derived below)
this gives $c_s^2\approx-0.84<0$, indicating gradient instability in the
\emph{pure} kinetic sector.  The geometric bulk-viscosity term modifies the
effective perturbation equations and can restore $c_s^2>0$; a full stability
analysis requiring the Israel--Stewart perturbation equations is deferred to
Paper~4.  The EFT embedding presented in this section should be regarded as
a \emph{motivational framework} for the algebraic structure of $w_0$, $w_a$,
not as a complete perturbation-stable theory.

\subsection{The SSEE Action}

The SSEE dark-energy sector is described by a k-essence scalar field $\phi$
minimally coupled to gravity and extended by a geometric bulk-viscosity term
in the Eckart approximation \citep{Eckart1940, Weinberg1971}:
%
\begin{equation}
  S = \int d^{4}x\,\sqrt{-g}
      \left[\frac{R}{16\pi G} + P(X) - \zeta\bigl(\nabla_{\!\mu}u^{\mu}\bigr)^{2}
            + \mathcal{L}_{m}\right],
  \label{eq:eft_action}
\end{equation}
%
where $X \equiv -\tfrac{1}{2}g^{\mu\nu}\partial_{\mu}\phi\,\partial_{\nu}\phi$
is the canonical kinetic term.  The fluid four-velocity of the dark-energy
component is identified with the normalised gradient of the field,
%
\begin{equation}
  u_{\mu} = -\frac{\partial_{\mu}\phi}{\sqrt{2X}},
  \label{eq:four_velocity}
\end{equation}
%
so that $u_{\mu}u^{\mu}=-1$ is satisfied automatically, and the viscous and
kinetic sectors remain self-consistently defined by the single degree of freedom
$\phi$.

% ─────────────────────────────────────────────────────────────────────────────
\subsection{Algebraic Determination of \texorpdfstring{$P(X)$}{P(X)}}
\label{sec:eft:PX}

For a power-law k-essence Lagrangian
$P(X) = \mathcal{A}\,X^{n}$ the equation of state is \emph{exactly} constant:
%
\begin{equation}
  w_{\phi}
    = \frac{P}{2X P_{X} - P}
    = \frac{\mathcal{A}X^{n}}{(2n-1)\mathcal{A}X^{n}}
    = \frac{1}{2n-1}.
  \label{eq:w_kessence}
\end{equation}
%
Imposing the SSEE algebraic constraint $w_{0} = -T_{r}/M_{v}$ uniquely fixes
the kinetic exponent without any fitting:
%
\begin{equation}
  \boxed{%
    n = \frac{T_{r} - M_{v}}{2\,T_{r}}
      = \frac{1 + w_{0}}{2w_{0}}
      \approx -0.09527,}
  \label{eq:n_ssee}
\end{equation}
%
where $T_{r} = 3(\varphi+\beta)$ and $M_{v} = \varphi+\pi+K_{v}$ are the
SSEE structural constants defined in Section~\ref{sec:axioms}.
One verifies immediately that $1/(2n-1) = -T_{r}/M_{v} = w_{0}$ identically.

The overall normalisation is fixed by requiring that the present-day
dark-energy density equals the SSEE prediction:
%
\begin{equation}
  \rho_{\mathrm{DE},0}
    = \frac{3H_{0}^{2}}{8\pi G}\,\frac{T_{r}}{M_{v}},
  \label{eq:rho_de0}
\end{equation}
%
with $H_{0}$ constrained observationally to
$66.75^{+0.44}_{-0.44}$~km\,s$^{-1}$\,Mpc$^{-1}$ by the MCMC analysis of
Paper~2.  Given the energy scale $\rho_{\mathrm{DE},0}$, the amplitude reads
%
\begin{equation}
  \mathcal{A}
    = \frac{\rho_{\mathrm{DE},0}^{1-n}}{2n-1},
  \label{eq:amplitude}
\end{equation}
%
which completes the specification of $P(X)$ with no free parameters beyond $H_{0}$.

\paragraph{Remark on the CPL evolution.}
The pure k-essence sector~(\ref{eq:w_kessence}) delivers a \emph{constant}
$w = w_{0}$.  The time variation $w_{a}$ arises from the viscous source
(Section~\ref{sec:eft:viscosity}), and is shown there to emerge algebraically
from the scale-factor evolution of the viscosity coefficient.

% ─────────────────────────────────────────────────────────────────────────────
\subsection{Geometric Bulk Viscosity and the Emergence of \texorpdfstring{$w_a$}{wa}}
\label{sec:eft:viscosity}

\subsubsection{Constant-coupling regime ($w_0$)}

The retention constant $\mathrm{KAL}_{0} \equiv \beta + \pi \approx 5.521$
governs the dissipative sector.  Following the Eckart first-order formalism,
the bulk viscosity coefficient is coupled to the instantaneous expansion rate:
%
\begin{equation}
  \zeta_0 = \frac{\mathrm{KAL}_{0}}{8\pi G}\,H.
  \label{eq:zeta}
\end{equation}
%
This delivers a \emph{constant-coupling} viscous pressure
$\Pi_0 = -3\zeta_0 H = -3\,\mathrm{KAL}_{0}H^{2}/8\pi G$,
which governs the background acceleration (Eq.~\ref{eq:Friedmann_a}) and fixes
the present-day effective equation of state $w_0 = -T_r/M_v$ algebraically
(Section~\ref{sec:eft:PX}).

\paragraph{Physical motivation for $\zeta \propto H$.}
Bulk viscosity proportional to the Hubble rate arises naturally in
kinetic theory when the dominant relaxation time is set by the expansion
itself \citep{Weinberg1971, Brevik2005}.  Within SSEE the coupling constant
is not fitted but predicted geometrically as
$\mathrm{KAL}_{0} = (\pi+\varphi)/2 + \pi$.

\subsubsection{Scale-factor evolution and the algebraic determination of \texorpdfstring{$w_a$}{wa}}
\label{sec:eft:wa_derivation}

The constant-coupling viscosity of the preceding section fixes $w_0$
algebraically but predicts a \emph{time-independent} effective equation of
state.  DESI~DR2 data require a time-varying component; within SSEE this must
arise from a scale-factor-dependent extension of the viscosity coefficient.

\paragraph{Structural requirements on $\zeta(a)$.}
Any admissible SSEE viscosity extension $\zeta(a)$ must satisfy three
independent constraints:
\begin{enumerate}[label=(\roman*)]
  \item \textbf{De~Sitter endpoint:} $\zeta(a)\to 0$ as $a\to 1$, so that no
        viscous term shifts the present-day equation of state away from the
        k-essence value $w_0 = -T_r/M_v$.
  \item \textbf{Dimensional consistency:} $[\zeta] = [\rho_{\mathrm{DE}}/H]$,
        constructed from the available dynamical quantities
        $\{\rho_{\mathrm{DE}},H,a\}$.
  \item \textbf{Algebraic closure:} the dimensionless amplitude must be a
        ratio of SSEE structural constants derived exclusively from
        $\{\varphi,\pi\}$.
\end{enumerate}
The unique admissible form satisfying (i) and (ii) at leading order in $(1-a)$
— the natural expansion around the de~Sitter endpoint — is
%
\begin{equation}
  \zeta(a) = \Gamma\,\frac{(1-a)\,\rho_{\mathrm{DE}}(a)}{H(a)},
  \label{eq:zeta_general}
\end{equation}
%
where $\Gamma$ is a dimensionless amplitude to be fixed by constraint (iii).

\paragraph{Algebraic determination of $\Gamma$.}
The k-essence sector has already employed the ratio $T_r/M_v$ to fix
$|w_0|$.  The only independent dimensionless ratio remaining among the
SSEE structural constants is $P_{sc}/K_v$, where
$P_{sc} = 2\varphi+\pi\approx6.378$ (Dynamical Evolution Scalar) and
$K_v = 2(\varphi+\pi)\approx9.519$ (Structural Constraint).
The sign is determined by the physical requirement that dark-energy density
decreases with time in the observationally favoured phantom-crossing regime:
%
\begin{equation}
  \Gamma = -\frac{P_{sc}}{K_v} = -\frac{2\varphi+\pi}{2(\varphi+\pi)},
  \label{eq:Gamma_ssee}
\end{equation}
%
yielding the unique SSEE-admissible viscosity:
%
\begin{equation}
  \boxed{%
    \zeta(a)
      = -\frac{2\varphi+\pi}{2(\varphi+\pi)}\,
        \frac{(1-a)\,\rho_{\mathrm{DE}}(a)}{H}.}
  \label{eq:zeta_ssee}
\end{equation}

\paragraph{Algebraic consistency with CPL.}
The CPL parametrisation $w(a) = w_0 + w_a(1-a)$ is the standard
model-independent description of slowly-varying dark energy; it serves here as
the observational language, not as a physical assumption of SSEE.  Adopting CPL
and substituting $\zeta(a)$ from Eq.~(\ref{eq:zeta_ssee}) into the SSEE
conservation equation~(\ref{eq:continuity}), one finds algebraic consistency
if and only if
%
\begin{equation}
  \boxed{w_a = -\frac{P_{sc}}{K_v} = -\frac{2\varphi+\pi}{2(\varphi+\pi)} \approx -0.6699.}
  \label{eq:wa_ssee}
\end{equation}
%
This is not a derived prediction in the sense that CPL is deduced from the
action; rather, it establishes that \emph{the SSEE viscosity uniquely fixes
the CPL coefficient $w_a$ from the algebraic structure of $\{\varphi,\pi\}$}.
The quantity $w_a$ is therefore not a free parameter of the model but a
computable consequence of the two-axiom system, in the same sense that $w_0$
is fixed by the k-essence exponent $n$.  At $a=1$ the viscous correction
vanishes ($\zeta\to 0$), recovering the pure k-essence regime $w = w_0$.

The implied dark-energy density evolution reads:
%
\begin{equation}
  \rho_{\mathrm{DE}}(a)
    = \rho_{\mathrm{DE},0}\,
      a^{-3(1+w_0+w_a)}\,e^{-3w_a(1-a)},
  \label{eq:rho_CPL}
\end{equation}
%
consistent with the CPL conservation equation
%
\begin{equation}
  \dot\rho_{\mathrm{DE}}
    + 3H\,\rho_{\mathrm{DE}}\,(1+w_0)
    = -3H\,w_a\,(1-a)\,\rho_{\mathrm{DE}},
  \label{eq:cont_CPL}
\end{equation}
%
with the effective viscous pressure
%
\begin{equation}
  \Pi(a)
    = w_a\,(1-a)\,\rho_{\mathrm{DE}}(a).
  \label{eq:Pi_evolving}
\end{equation}
%
At $a \ll 1$ the factor $(1-a)\to 1$ and $\rho_{\mathrm{DE}}$ redshifts
according to Eq.~(\ref{eq:rho_CPL}), so $\zeta$ remains finite and
well-behaved throughout cosmic history.

\paragraph{Boundary conditions.}
The full viscosity interpolates between two algebraically-fixed regimes:
%
\begin{align}
  a \to 0 \;(\text{early}): &\quad
    \zeta \to -\frac{2\varphi+\pi}{2(\varphi+\pi)}\,
              \frac{\rho_{\mathrm{DE,0}}\,a^{-3(1+w_0+w_a)}\,e^{-3w_a}}{H(a)},
  \label{eq:zeta_early}\\[4pt]
  a \to 1 \;(\text{today}): &\quad
    \zeta \to 0.
  \label{eq:zeta_today}
\end{align}
%
Both limits are free of divergences and require no new parameters.

The viscous pressure is therefore
%
\begin{equation}
  \Pi(a)
    = \underbrace{-\frac{3\,\mathrm{KAL}_{0}\,H^{2}}{8\pi G}}_{\text{background}}
      \underbrace{-\;3H\,w_a\,(1-a)\,\rho_{\mathrm{DE}}(a)}_{\text{CPL evolution}},
  \label{eq:Pi_total}
\end{equation}
%
where the first term is the constant-coupling Eckart pressure that sustains
$\ddot a > 0$ and the second term generates the running $w(a)$ observed by
DESI~DR2 \citep{AbdulKarim2025}.

% ─────────────────────────────────────────────────────────────────────────────
\subsection{Modified Friedmann Equations}
\label{sec:eft:friedmann}

Applying the principle of least action ($\delta S = 0$) to the
FLRW metric $ds^{2} = -dt^{2} + a^{2}(t)\,d\mathbf{x}^{2}$ yields
three governing equations for the SSEE background.

\paragraph{I. Energy equation:}
\begin{equation}
  H^{2}
    = \frac{8\pi G}{3}
      \bigl(\rho_{m} + \rho_{r} + \rho_{\mathrm{SSEE}}\bigr).
  \label{eq:Friedmann_H}
\end{equation}

\paragraph{II. Acceleration equation:}
\begin{equation}
  \frac{\ddot{a}}{a}
    = -\frac{4\pi G}{3}
      \!\left[\rho_{m} + 2\rho_{r}
        + \rho_{\mathrm{SSEE}}\!\left(1 + 3w_{0}\right)
        - \frac{9\,\mathrm{KAL}_{0}\,H^{2}}{8\pi G}
      \right].
  \label{eq:Friedmann_a}
\end{equation}

\paragraph{III. Conservation equation with dissipative source:}
\begin{equation}
  \dot{\rho}_{\mathrm{SSEE}}
    + 3H\,\rho_{\mathrm{SSEE}}\!\left(1 + w_{0}\right)
    = \frac{9\,\mathrm{KAL}_{0}\,H^{3}}{8\pi G}.
  \label{eq:continuity}
\end{equation}

\noindent\textbf{Derivation of the coefficient 9.}\par\smallskip
\noindent From Eq.~(\ref{eq:zeta}), the viscous pressure is
$\Pi = -3\mathrm{KAL}_{0}H^{2}/8\pi G$.
In the Raychaudhuri equation the effective pressure of the SSEE sector enters as
$p_{\mathrm{eff}} = w_{0}\rho_{\mathrm{SSEE}} + \Pi$, so the acceleration term becomes
%
\begin{align}
  -\frac{4\pi G}{3}\bigl(\rho_{\mathrm{SSEE}} + 3p_{\mathrm{eff}}\bigr)
  &= -\frac{4\pi G}{3}
     \!\left[\rho_{\mathrm{SSEE}}(1+3w_{0}) + 3\Pi\right] \notag\\
  &= -\frac{4\pi G}{3}
     \!\left[\rho_{\mathrm{SSEE}}(1+3w_{0})
             - \frac{9\,\mathrm{KAL}_{0}\,H^{2}}{8\pi G}\right],
  \label{eq:deriv_coeff9}
\end{align}
which is exactly Eq.~(\ref{eq:Friedmann_a}).  The factor 9 arises from
$3 \times |\Pi|/(H^{2}) = 3 \times 3\mathrm{KAL}_{0}/8\pi G$; a factor
of 3 rather than 9 would correspond to including $\Pi$ only once instead
of as the trace of the viscous stress.

Similarly, the conservation law $\nabla_{\!\mu}T^{\mu\nu}=0$ for the effective
fluid $(w_{0},\Pi)$ gives
$\dot{\rho} + 3H(\rho + w_{0}\rho + \Pi) = 0$, hence
$\dot{\rho} + 3H\rho(1+w_{0}) = 9\mathrm{KAL}_{0}H^{3}/8\pi G$,
confirming that Eq.~(\ref{eq:Friedmann_a}) and Eq.~(\ref{eq:continuity})
are mutually consistent.

% ─────────────────────────────────────────────────────────────────────────────
\subsection{Israel--Stewart Viscous Perturbations}
\label{sec:eft:IS}

The Eckart formulation used in \S\ref{sec:eft:PX}--\ref{sec:eft:friedmann}
is a first-order theory that violates causality \citep{HiscockLindblom1985}:
the viscous signal propagates instantaneously.  Israel--Stewart (IS) theory
\citep{IsraelStewart1979} resolves this by promoting $\Pi$ to a dynamical
variable with a finite relaxation time $\tau_{\Pi}$:
%
\begin{equation}
  \tau_{\Pi} H_{0}\,\tilde{h}(a)\,
    \frac{\mathrm{d}\tilde{\Pi}}{\mathrm{d}\ln a}
  + \tilde{\Pi}
  = -\,\mathrm{KAL}_{0}\,\tilde{h}^{2}(a),
  \label{eq:IS_transport}
\end{equation}
%
where $\tilde{h} = H/H_{0}$, $\tilde{\Pi} = \Pi\,8\pi G/H_{0}^{2}$, and
all densities are normalised to $\rho_{\mathrm{crit},0}$.  In the
limit $\tau_{\Pi} \to 0$, Eq.~(\ref{eq:IS_transport}) reduces to the
Eckart result $\tilde{\Pi} = -\mathrm{KAL}_{0}\tilde{h}^{2}$.

\paragraph{Causality constraint.}
The bulk-viscous sound speed in IS theory is \citep{HiscockLindblom1985}
%
\begin{equation}
  c_{s}^{2}
    = \frac{\zeta}{\rho_{\mathrm{DE}}\,\tau_{\Pi}}
    = \frac{\mathrm{KAL}_{0}}{3\,\Omega_{\mathrm{DE}}\,(\tau_{\Pi} H_{0})}
    \leq 1.
  \label{eq:IS_causality}
\end{equation}
%
Setting $c_{s}^{2} = 1$ (the causality boundary) uniquely fixes the
IS relaxation time as an algebraic combination of SSEE constants:
%
\begin{equation}
  \tau_{\Pi} H_{0}
    \;=\; \frac{\mathrm{KAL}_{0}}{3\,\Omega_{\mathrm{DE}}}
    \;=\; \frac{\mathrm{KAL}_{0}\,M_{v}}{3\,T_{r}}
    \;\approx\; 2.191,
  \label{eq:tau_IS}
\end{equation}
%
with no free parameters.

\paragraph{Growth index.}
Coupling Eq.~(\ref{eq:IS_transport}) to the conservation
equation~(\ref{eq:continuity}) and the linear-growth ODE
%
\begin{equation}
  \delta'' + \left[2 + \frac{\mathrm{d}\ln\tilde{h}}{\mathrm{d}\ln a}\right]\delta'
  = \frac{3}{2}\,\frac{\Omega_{m,\mathrm{dyn}}\,a^{-3}}{\tilde{h}^{2}(a)}\,\delta,
  \label{eq:growth_IS}
\end{equation}
%
we solve the system numerically (\texttt{ssee\_is\_growth.py}) with
the IS background $\tilde{h}^{2}(a) = \Omega_{m,\mathrm{dyn}}\,a^{-3}+\rho_{\mathrm{DE}}(a)$.
The growth rate $f \equiv \mathrm{d}\ln\delta/\mathrm{d}\ln a \approx \Omega_{m}(a)^{\gamma}$
is well described by a power law over $0.1 \leq a \leq 1$ with
%
\begin{equation}
  \gamma_{\mathrm{IS}} = 0.657 \pm 0.002
  \quad(\text{essentially independent of }\tau_{\Pi} H_{0}).
  \label{eq:gamma_IS}
\end{equation}
%
This differs by 6\% from the algebraic sovereignty $\gamma = \varphi^{-1} = 0.618$.
Paper~3 \S5.4 reports both values: the algebraic prediction $\gamma_\mathrm{alg}=0.618$
and the IS numerical result $\gamma_\mathrm{IS}=0.657$.
The IS formalism provides the causal theoretical structure; whether $\gamma = \varphi^{-1}$
is the true attractor of the IS boundary conditions remains a target prediction
for forthcoming Stage-IV weak-lensing surveys (LSST, Euclid).
Note that the direct IS growth-factor integral (this Appendix) yields $S_8=0.837$,
while Paper~3's CAMB post-processing method with the same $\gamma_\mathrm{IS}=0.657$
gives $S_8=0.818$; both are consistent with Planck 2018 ($0.832\pm0.013$) within $1\sigma$.

\paragraph{$S_{8}$ constraint.}
Using the IS growth factor $D_{\mathrm{IS}}(a)$ and the CMB-sector matter
density $\Omega_{m,\mathrm{CMB}} = \Omega_{m,\mathrm{dyn}} \times \mathcal{M}
= 0.3199$ (MIRA, Paper~3),
%
\begin{equation}
  S_{8}^{\mathrm{IS}}
    \equiv \sigma_{8,\mathrm{IS}}
           \sqrt{\frac{\Omega_{m,\mathrm{CMB}}}{0.3}}
    \approx 0.837,
  \label{eq:S8_IS}
\end{equation}
%
within $0.4\sigma$ of the Planck-2018 value $0.832 \pm 0.013$~\citep{PlanckVI2020}.
Using the dynamic sector $\Omega_{m,\mathrm{dyn}} = 0.160$ instead gives
$S_{8} \approx 0.59$ ($6\sigma$ tension); the correct physical choice depends
on the formal resolution of the two-sector Friedmann problem (Section~\ref{sec:eft:status}, item 3).

% ─────────────────────────────────────────────────────────────────────────────
\subsection{Physical Regimes}
\label{sec:eft:regimes}

\paragraph{Late universe ($z \lesssim 2$, DESI~DR2).}
At low redshift $H \approx H_{0}$, so the viscous source
$\propto \mathrm{KAL}_{0}H^{3}/8\pi G$ is small compared to
$\rho_{\mathrm{SSEE}}$.  The acceleration equation is dominated by
$\rho_{\mathrm{SSEE}}(1+3w_{0}) < 0$
(since $w_{0} \approx -0.840 < -1/3$), driving $\ddot{a} > 0$
without dark matter.  This regime is the one tested against DESI~DR2
in Paper~2, where the algebraic prediction $(w_{0},w_{a})$ matches
the data at $0.05\sigma$.

\paragraph{Early universe (CMB epoch, $z \sim 10^{3}$).}
At recombination $H \gg H_{0}$, the geometric retention term
$9\,\mathrm{KAL}_{0}H^{3}/8\pi G$ acts as a dissipative source in
Eq.~(\ref{eq:continuity}), injecting energy into the dark-energy fluid
and modifying its effective equation of state.  This structural
viscosity provides the physical mechanism by which an effective
matter density $\Omega_{m,\mathrm{eff}} = 1 - T_{r}/M_{v} \approx 0.160$
yields the correct CMB acoustic scale through the MIRA factor
$\mathcal{M} = 2 - 1/\varphi^{2} \approx 1.9989$ (Paper~3, \S2.3).

% ─────────────────────────────────────────────────────────────────────────────
% ─────────────────────────────────────────────────────────────────────────────
\subsection{Inflationary Embedding: \texorpdfstring{$\alpha$}{α}-Attractor with \texorpdfstring{$\alpha=\varphi^{4}/3$}{α=φ⁴/3}}
\label{sec:eft:inflation}

The SSEE algebraic predictions for the primordial scalar spectral index
$n_s = 1 - \varphi^{-7}$ (Paper~1, \S\ref{sec:register}) and the
tensor-to-scalar ratio $r = \varphi^{-10}$ (Paper~4) belong simultaneously
to a well-defined class of inflationary models: the
\emph{$\alpha$-attractors} of Kallosh \& Linde \citeyearpar{KalloshLinde2013}.

\paragraph{$\alpha$-attractor predictions.}
In the universal leading-order limit, $\alpha$-attractor models predict
\begin{equation}
  n_s = 1 - \frac{2}{N}, \qquad
  r   = \frac{12\alpha}{N^{2}},
  \label{eq:alpha_attractor}
\end{equation}
where $N$ is the number of inflationary e-folds and $\alpha > 0$
parameterises the inverse curvature of the scalar field's K\"{a}hler target
manifold \citep{KalloshLinde2013,Ferrara2013}.
Setting $\alpha = 1$ recovers Starobinsky $R^{2}$ inflation exactly.

\paragraph{SSEE identification.}
Inserting $n_s = 1 - \varphi^{-7}$ into the first relation gives
\begin{equation}
  N = 2/\varphi^{-7} = 2\varphi^{7} \approx 58.07,
  \label{eq:N_ssee}
\end{equation}
which lies squarely in the observationally required window $50\lesssim N\lesssim 60$.
Substituting $N = 2\varphi^{7}$ and $r = \varphi^{-10}$ into the second relation:
\begin{equation}
  \alpha = \frac{r\,N^{2}}{12}
         = \frac{\varphi^{-10}\cdot 4\varphi^{14}}{12}
         = \frac{4\varphi^{4}}{12}
         = \frac{\varphi^{4}}{3}.
  \label{eq:alpha_ssee}
\end{equation}
This result is algebraically exact (verified numerically to machine precision;
see \texttt{ssee\_inflation\_connection.py}).
Using the Fibonacci identity $\varphi^{4} = 3\varphi + 2$:
\begin{equation}
  \alpha = \frac{3\varphi + 2}{3} = \varphi + \frac{2}{3} \approx 2.285.
  \label{eq:alpha_fibonacci}
\end{equation}

\paragraph{Geometric interpretation.}
In the $\alpha$-attractor framework the K\"{a}hler curvature of the
inflaton's target space (a Poincar\'{e} disk of radius $\sqrt{3\alpha}$)
is $\mathcal{R} = -2/(3\alpha)$.
For $\alpha = \varphi^{4}/3$:
\begin{equation}
  \mathcal{R} = -\frac{2}{3 \cdot \varphi^{4}/3} = -\frac{2}{\varphi^{4}}
               = -\varphi^{-4} \approx -0.146.
  \label{eq:kahler_curvature}
\end{equation}
The inflaton therefore lives on a hyperbolic manifold whose curvature is
set by $\varphi$ alone. This gives a \emph{geometric} origin for the
appearance of $\varphi$ in inflationary observables: $\varphi$ enters not
as a numerical coincidence but as the intrinsic length scale of the
K\"{a}hler geometry of the SSEE inflaton sector.

\paragraph{Slow-roll parameters.}
The SSEE predictions $(n_s,\,r) = (1-\varphi^{-7},\,\varphi^{-10})$ fix the
slow-roll parameters uniquely:
\begin{align}
  \varepsilon &= \frac{r}{16} = \frac{\varphi^{-10}}{16}, \label{eq:sr_eps} \\
  \eta &= \frac{n_s - 1 + 6\varepsilon}{2}
        = \frac{-\varphi^{-7} + 3\varphi^{-10}/8}{2}, \label{eq:sr_eta} \\
  \frac{\eta}{\varepsilon} &= 3 - 8\varphi^{3} = -5 - 16\varphi \approx -30.9.
  \label{eq:sr_ratio}
\end{align}
The ratio $\eta/\varepsilon = 3 - 8\varphi^{3}$ is a pure algebraic
identity derived from the SSEE axioms $\{\varphi, \pi\}$; it defines a
one-parameter family of inflationary potentials $V(\phi_{\rm inf})$ whose
single element consistent with $\alpha$-attractors is $\alpha = \varphi^{4}/3$.

\paragraph{Relation to Starobinsky.}
Starobinsky ($\alpha = 1$) reproduces $n_s$ correctly for $N = 2\varphi^{7}$
but gives $r_{\rm Staro} = 3/\varphi^{14} \approx 0.00356$, a factor of
$\varphi^{4}/3\approx 2.28$ below the SSEE prediction $\varphi^{-10}$.
SSEE is therefore \emph{not} Starobinsky inflation; it is a generalised
$\alpha$-attractor that reduces to Starobinsky only in the limit
$\varphi \to 1$ (i.e.\ $\alpha \to 1$, the fixed point of the golden-ratio
recursion $x = 1 + 1/x$).

% ─────────────────────────────────────────────────────────────────────────────
\subsection{Summary and Open Problems}
\label{sec:eft:status}

Table~\ref{tab:eft_summary} summarises the EFT status.
The action~(\ref{eq:eft_action}) with kinetic exponent~(\ref{eq:n_ssee})
and viscosity~(\ref{eq:zeta_ssee}) constitutes a \emph{complete},
zero-free-parameter EFT description of the SSEE background evolution.
Every constant ($n$, $\mathrm{KAL}_{0}$, $w_0$, $w_a$, $\Omega_{\mathrm{DE}}$) is
derived from $\varphi$ and $\pi$ alone; only the overall energy scale $H_{0}$
enters as an observational input.  Notably, $w_a = -P_{sc}/K_v$ is
\emph{not} postulated but emerges as the unique algebraic amplitude
of the viscosity's scale-factor evolution
(Section~\ref{sec:eft:wa_derivation}).

Four aspects require further development:
%
\begin{enumerate}
  \item \textbf{Growth index.}
        The IS-derived growth index $\gamma_{\mathrm{IS}} = 0.657$
        (\S\ref{sec:eft:IS}) differs by 6\% from the algebraic prediction
        $\varphi^{-1} = 0.618$.  A mechanism that would select
        $\gamma = \varphi^{-1}$ from IS boundary conditions is not yet
        identified; this constitutes a near-prediction testable by
        LSST/Euclid Stage-IV weak lensing.

  \item \textbf{Full perturbation spectrum.}
        The IS growth ODE~(\ref{eq:growth_IS}) governs sub-Hubble dark-matter
        perturbations; the matching to CMB scalar transfer functions and
        the tensor-mode propagation in the IS background require a complete
        Boltzmann implementation.

  \item \textbf{Radiative stability.}
        The EFT cutoff scale $\Lambda_{\mathrm{UV}}$ and the quantum stability
        of the $n < 0$ power-law sector are not yet established.

  \item \textbf{Inflation--dark energy unification.}
        The $\alpha$-attractor sector (\S\ref{sec:eft:inflation}) and the
        k-essence dark-energy sector (\S\ref{sec:eft:PX}) both originate
        from $\varphi$, suggesting a single unified action.
        Constructing the quintessential inflation potential $V(\phi)$
        that reduces to $P(X) = \mathcal{A}X^n$ at late times
        and to an $\alpha=\varphi^4/3$ attractor at early times
        is identified as high-priority future work.
\end{enumerate}

\begin{table}[h]
\centering
\caption{EFT status summary for SSEE-V3.6.}
\label{tab:eft_summary}
\begin{tabular}{lll}
\hline\hline
Element & Status & Derivation \\
\hline
Action form $S$ & \checkmark Complete & Standard GR + k-essence + Eckart \\
$u_{\mu}(\phi)$ coupling & \checkmark Complete & $u_{\mu} = -\partial_{\mu}\phi/\sqrt{2X}$ \\
Kinetic exponent $n$ & \checkmark Algebraic & $n=(T_{r}-M_{v})/2T_{r}$ from $\varphi,\pi$ \\
Equation of state $w_{0}$ & \checkmark Algebraic & $w_{0}=1/(2n-1)=-T_{r}/M_{v}$ \\
Viscosity coupling $\zeta$ & \checkmark Algebraic & $\zeta=\mathrm{KAL}_{0}H/8\pi G$ \\
Friedmann eqs.\ I--III & \checkmark Consistent & Factor 9 verified in \S\ref{sec:eft:friedmann} \\
$w_{a}$ from EFT & \checkmark Algebraic & $\zeta(a)\!\propto\!(1{-}a)\rho_\mathrm{DE}/H$; see \S\ref{sec:eft:wa_derivation} \\
IS relaxation time $\tau_\Pi$ & \checkmark Algebraic & $\tau_\Pi H_0 = \mathrm{KAL}_0/(3\Omega_\mathrm{DE}) \approx 2.191$ \\
IS causality ($c_s^2 \leq 1$) & \checkmark Verified & Eq.~(\ref{eq:IS_causality}); $c_s^2 = 1$ at boundary \\
Growth index $\gamma_\mathrm{IS}$ & \checkmark Numerical & $0.657 \pm 0.002$ (6\% above $\varphi^{-1}=0.618$) \\
$S_8$ (CMB sector) & \checkmark Numerical & $0.837$, within $0.4\sigma$ of Planck \\
$\alpha$-attractor ($\alpha=\varphi^4/3$) & \checkmark Algebraic & Eq.~(\ref{eq:alpha_ssee}); exact, $N=2\varphi^7$ \\
K\"{a}hler curvature $\mathcal{R}=-\varphi^{-4}$ & \checkmark Algebraic & Eq.~(\ref{eq:kahler_curvature}) \\
Full perturbation spectrum & $\circ$ Open & Off-diagonal CMB cov.\ + Stage-IV WL \\
Inflation--DE unification & $\circ$ Open & Quintessential inflation potential \\
Radiative stability & $\circ$ Open & Future work \\
\hline
\end{tabular}
\end{table}

%% All equations verified algebraically (see src/ssee_eft_verify.py)

\section{Effective Field Theory Embedding}
\label{sec:eft}

\textbf{Scope and limitations of this section.}
The background dynamics presented here use the \emph{Eckart first-order}
approximation \citep{Eckart1940}, which is sufficient to derive $w_0$, $w_a$,
and the viscosity coupling $\mathrm{KAL}_0$ on superhorizon scales.
The Eckart formulation is known to violate causality in the relativistic
perturbation sector \citep{HiscockLindblom1985}; for perturbation-level predictions
(B-mode forecasts, sound-speed constraints) the Israel--Stewart (IS) formulation
is required and is developed in Paper~4 \S\ref{sec:is_perturbations}.
The Eckart and IS frameworks agree at leading order in the slow-roll expansion,
so all background results below are unchanged by the IS extension.

\paragraph{Sound speed caveat (EFT embedding).}
For a pure power-law k-essence Lagrangian $P(X)=A\,X^n$ the adiabatic
sound speed satisfies $c_s^2 = 1/(2n-1)$.  With $n\approx-0.095$ (derived below)
this gives $c_s^2\approx-0.84<0$, indicating gradient instability in the
\emph{pure} kinetic sector.  The geometric bulk-viscosity term modifies the
effective perturbation equations and can restore $c_s^2>0$; a full stability
analysis requiring the Israel--Stewart perturbation equations is deferred to
Paper~4.  The EFT embedding presented in this section should be regarded as
a \emph{motivational framework} for the algebraic structure of $w_0$, $w_a$,
not as a complete perturbation-stable theory.

\subsection{The SSEE Action}

The SSEE dark-energy sector is described by a k-essence scalar field $\phi$
minimally coupled to gravity and extended by a geometric bulk-viscosity term
in the Eckart approximation \citep{Eckart1940, Weinberg1971}:
%
\begin{equation}
  S = \int d^{4}x\,\sqrt{-g}
      \left[\frac{R}{16\pi G} + P(X) - \zeta\bigl(\nabla_{\!\mu}u^{\mu}\bigr)^{2}
            + \mathcal{L}_{m}\right],
  \label{eq:eft_action}
\end{equation}
%
where $X \equiv -\tfrac{1}{2}g^{\mu\nu}\partial_{\mu}\phi\,\partial_{\nu}\phi$
is the canonical kinetic term.  The fluid four-velocity of the dark-energy
component is identified with the normalised gradient of the field,
%
\begin{equation}
  u_{\mu} = -\frac{\partial_{\mu}\phi}{\sqrt{2X}},
  \label{eq:four_velocity}
\end{equation}
%
so that $u_{\mu}u^{\mu}=-1$ is satisfied automatically, and the viscous and
kinetic sectors remain self-consistently defined by the single degree of freedom
$\phi$.

% ─────────────────────────────────────────────────────────────────────────────
\subsection{Algebraic Determination of \texorpdfstring{$P(X)$}{P(X)}}
\label{sec:eft:PX}

For a power-law k-essence Lagrangian
$P(X) = \mathcal{A}\,X^{n}$ the equation of state is \emph{exactly} constant:
%
\begin{equation}
  w_{\phi}
    = \frac{P}{2X P_{X} - P}
    = \frac{\mathcal{A}X^{n}}{(2n-1)\mathcal{A}X^{n}}
    = \frac{1}{2n-1}.
  \label{eq:w_kessence}
\end{equation}
%
Imposing the SSEE algebraic constraint $w_{0} = -T_{r}/M_{v}$ uniquely fixes
the kinetic exponent without any fitting:
%
\begin{equation}
  \boxed{%
    n = \frac{T_{r} - M_{v}}{2\,T_{r}}
      = \frac{1 + w_{0}}{2w_{0}}
      \approx -0.09527,}
  \label{eq:n_ssee}
\end{equation}
%
where $T_{r} = 3(\varphi+\beta)$ and $M_{v} = \varphi+\pi+K_{v}$ are the
SSEE structural constants defined in Section~\ref{sec:axioms}.
One verifies immediately that $1/(2n-1) = -T_{r}/M_{v} = w_{0}$ identically.

The overall normalisation is fixed by requiring that the present-day
dark-energy density equals the SSEE prediction:
%
\begin{equation}
  \rho_{\mathrm{DE},0}
    = \frac{3H_{0}^{2}}{8\pi G}\,\frac{T_{r}}{M_{v}},
  \label{eq:rho_de0}
\end{equation}
%
with $H_{0}$ constrained observationally to
$66.75^{+0.44}_{-0.44}$~km\,s$^{-1}$\,Mpc$^{-1}$ by the MCMC analysis of
Paper~2.  Given the energy scale $\rho_{\mathrm{DE},0}$, the amplitude reads
%
\begin{equation}
  \mathcal{A}
    = \frac{\rho_{\mathrm{DE},0}^{1-n}}{2n-1},
  \label{eq:amplitude}
\end{equation}
%
which completes the specification of $P(X)$ with no free parameters beyond $H_{0}$.

\paragraph{Remark on the CPL evolution.}
The pure k-essence sector~(\ref{eq:w_kessence}) delivers a \emph{constant}
$w = w_{0}$.  The time variation $w_{a}$ arises from the viscous source
(Section~\ref{sec:eft:viscosity}), and is shown there to emerge algebraically
from the scale-factor evolution of the viscosity coefficient.

% ─────────────────────────────────────────────────────────────────────────────
\subsection{Geometric Bulk Viscosity and the Emergence of \texorpdfstring{$w_a$}{wa}}
\label{sec:eft:viscosity}

\subsubsection{Constant-coupling regime ($w_0$)}

The retention constant $\mathrm{KAL}_{0} \equiv \beta + \pi \approx 5.521$
governs the dissipative sector.  Following the Eckart first-order formalism,
the bulk viscosity coefficient is coupled to the instantaneous expansion rate:
%
\begin{equation}
  \zeta_0 = \frac{\mathrm{KAL}_{0}}{8\pi G}\,H.
  \label{eq:zeta}
\end{equation}
%
This delivers a \emph{constant-coupling} viscous pressure
$\Pi_0 = -3\zeta_0 H = -3\,\mathrm{KAL}_{0}H^{2}/8\pi G$,
which governs the background acceleration (Eq.~\ref{eq:Friedmann_a}) and fixes
the present-day effective equation of state $w_0 = -T_r/M_v$ algebraically
(Section~\ref{sec:eft:PX}).

\paragraph{Physical motivation for $\zeta \propto H$.}
Bulk viscosity proportional to the Hubble rate arises naturally in
kinetic theory when the dominant relaxation time is set by the expansion
itself \citep{Weinberg1971, Brevik2005}.  Within SSEE the coupling constant
is not fitted but predicted geometrically as
$\mathrm{KAL}_{0} = (\pi+\varphi)/2 + \pi$.

\subsubsection{Scale-factor evolution and the algebraic determination of \texorpdfstring{$w_a$}{wa}}
\label{sec:eft:wa_derivation}

The constant-coupling viscosity of the preceding section fixes $w_0$
algebraically but predicts a \emph{time-independent} effective equation of
state.  DESI~DR2 data require a time-varying component; within SSEE this must
arise from a scale-factor-dependent extension of the viscosity coefficient.

\paragraph{Structural requirements on $\zeta(a)$.}
Any admissible SSEE viscosity extension $\zeta(a)$ must satisfy three
independent constraints:
\begin{enumerate}[label=(\roman*)]
  \item \textbf{De~Sitter endpoint:} $\zeta(a)\to 0$ as $a\to 1$, so that no
        viscous term shifts the present-day equation of state away from the
        k-essence value $w_0 = -T_r/M_v$.
  \item \textbf{Dimensional consistency:} $[\zeta] = [\rho_{\mathrm{DE}}/H]$,
        constructed from the available dynamical quantities
        $\{\rho_{\mathrm{DE}},H,a\}$.
  \item \textbf{Algebraic closure:} the dimensionless amplitude must be a
        ratio of SSEE structural constants derived exclusively from
        $\{\varphi,\pi\}$.
\end{enumerate}
The unique admissible form satisfying (i) and (ii) at leading order in $(1-a)$
— the natural expansion around the de~Sitter endpoint — is
%
\begin{equation}
  \zeta(a) = \Gamma\,\frac{(1-a)\,\rho_{\mathrm{DE}}(a)}{H(a)},
  \label{eq:zeta_general}
\end{equation}
%
where $\Gamma$ is a dimensionless amplitude to be fixed by constraint (iii).

\paragraph{Algebraic determination of $\Gamma$.}
The k-essence sector has already employed the ratio $T_r/M_v$ to fix
$|w_0|$.  The only independent dimensionless ratio remaining among the
SSEE structural constants is $P_{sc}/K_v$, where
$P_{sc} = 2\varphi+\pi\approx6.378$ (Dynamical Evolution Scalar) and
$K_v = 2(\varphi+\pi)\approx9.519$ (Structural Constraint).
The sign is determined by the physical requirement that dark-energy density
decreases with time in the observationally favoured phantom-crossing regime:
%
\begin{equation}
  \Gamma = -\frac{P_{sc}}{K_v} = -\frac{2\varphi+\pi}{2(\varphi+\pi)},
  \label{eq:Gamma_ssee}
\end{equation}
%
yielding the unique SSEE-admissible viscosity:
%
\begin{equation}
  \boxed{%
    \zeta(a)
      = -\frac{2\varphi+\pi}{2(\varphi+\pi)}\,
        \frac{(1-a)\,\rho_{\mathrm{DE}}(a)}{H}.}
  \label{eq:zeta_ssee}
\end{equation}

\paragraph{Algebraic consistency with CPL.}
The CPL parametrisation $w(a) = w_0 + w_a(1-a)$ is the standard
model-independent description of slowly-varying dark energy; it serves here as
the observational language, not as a physical assumption of SSEE.  Adopting CPL
and substituting $\zeta(a)$ from Eq.~(\ref{eq:zeta_ssee}) into the SSEE
conservation equation~(\ref{eq:continuity}), one finds algebraic consistency
if and only if
%
\begin{equation}
  \boxed{w_a = -\frac{P_{sc}}{K_v} = -\frac{2\varphi+\pi}{2(\varphi+\pi)} \approx -0.6699.}
  \label{eq:wa_ssee}
\end{equation}
%
This is not a derived prediction in the sense that CPL is deduced from the
action; rather, it establishes that \emph{the SSEE viscosity uniquely fixes
the CPL coefficient $w_a$ from the algebraic structure of $\{\varphi,\pi\}$}.
The quantity $w_a$ is therefore not a free parameter of the model but a
computable consequence of the two-axiom system, in the same sense that $w_0$
is fixed by the k-essence exponent $n$.  At $a=1$ the viscous correction
vanishes ($\zeta\to 0$), recovering the pure k-essence regime $w = w_0$.

The implied dark-energy density evolution reads:
%
\begin{equation}
  \rho_{\mathrm{DE}}(a)
    = \rho_{\mathrm{DE},0}\,
      a^{-3(1+w_0+w_a)}\,e^{-3w_a(1-a)},
  \label{eq:rho_CPL}
\end{equation}
%
consistent with the CPL conservation equation
%
\begin{equation}
  \dot\rho_{\mathrm{DE}}
    + 3H\,\rho_{\mathrm{DE}}\,(1+w_0)
    = -3H\,w_a\,(1-a)\,\rho_{\mathrm{DE}},
  \label{eq:cont_CPL}
\end{equation}
%
with the effective viscous pressure
%
\begin{equation}
  \Pi(a)
    = w_a\,(1-a)\,\rho_{\mathrm{DE}}(a).
  \label{eq:Pi_evolving}
\end{equation}
%
At $a \ll 1$ the factor $(1-a)\to 1$ and $\rho_{\mathrm{DE}}$ redshifts
according to Eq.~(\ref{eq:rho_CPL}), so $\zeta$ remains finite and
well-behaved throughout cosmic history.

\paragraph{Boundary conditions.}
The full viscosity interpolates between two algebraically-fixed regimes:
%
\begin{align}
  a \to 0 \;(\text{early}): &\quad
    \zeta \to -\frac{2\varphi+\pi}{2(\varphi+\pi)}\,
              \frac{\rho_{\mathrm{DE,0}}\,a^{-3(1+w_0+w_a)}\,e^{-3w_a}}{H(a)},
  \label{eq:zeta_early}\\[4pt]
  a \to 1 \;(\text{today}): &\quad
    \zeta \to 0.
  \label{eq:zeta_today}
\end{align}
%
Both limits are free of divergences and require no new parameters.

The viscous pressure is therefore
%
\begin{equation}
  \Pi(a)
    = \underbrace{-\frac{3\,\mathrm{KAL}_{0}\,H^{2}}{8\pi G}}_{\text{background}}
      \underbrace{-\;3H\,w_a\,(1-a)\,\rho_{\mathrm{DE}}(a)}_{\text{CPL evolution}},
  \label{eq:Pi_total}
\end{equation}
%
where the first term is the constant-coupling Eckart pressure that sustains
$\ddot a > 0$ and the second term generates the running $w(a)$ observed by
DESI~DR2 \citep{AbdulKarim2025}.

% ─────────────────────────────────────────────────────────────────────────────
\subsection{Modified Friedmann Equations}
\label{sec:eft:friedmann}

Applying the principle of least action ($\delta S = 0$) to the
FLRW metric $ds^{2} = -dt^{2} + a^{2}(t)\,d\mathbf{x}^{2}$ yields
three governing equations for the SSEE background.

\paragraph{I. Energy equation:}
\begin{equation}
  H^{2}
    = \frac{8\pi G}{3}
      \bigl(\rho_{m} + \rho_{r} + \rho_{\mathrm{SSEE}}\bigr).
  \label{eq:Friedmann_H}
\end{equation}

\paragraph{II. Acceleration equation:}
\begin{equation}
  \frac{\ddot{a}}{a}
    = -\frac{4\pi G}{3}
      \!\left[\rho_{m} + 2\rho_{r}
        + \rho_{\mathrm{SSEE}}\!\left(1 + 3w_{0}\right)
        - \frac{9\,\mathrm{KAL}_{0}\,H^{2}}{8\pi G}
      \right].
  \label{eq:Friedmann_a}
\end{equation}

\paragraph{III. Conservation equation with dissipative source:}
\begin{equation}
  \dot{\rho}_{\mathrm{SSEE}}
    + 3H\,\rho_{\mathrm{SSEE}}\!\left(1 + w_{0}\right)
    = \frac{9\,\mathrm{KAL}_{0}\,H^{3}}{8\pi G}.
  \label{eq:continuity}
\end{equation}

\noindent\textbf{Derivation of the coefficient 9.}\par\smallskip
\noindent From Eq.~(\ref{eq:zeta}), the viscous pressure is
$\Pi = -3\mathrm{KAL}_{0}H^{2}/8\pi G$.
In the Raychaudhuri equation the effective pressure of the SSEE sector enters as
$p_{\mathrm{eff}} = w_{0}\rho_{\mathrm{SSEE}} + \Pi$, so the acceleration term becomes
%
\begin{align}
  -\frac{4\pi G}{3}\bigl(\rho_{\mathrm{SSEE}} + 3p_{\mathrm{eff}}\bigr)
  &= -\frac{4\pi G}{3}
     \!\left[\rho_{\mathrm{SSEE}}(1+3w_{0}) + 3\Pi\right] \notag\\
  &= -\frac{4\pi G}{3}
     \!\left[\rho_{\mathrm{SSEE}}(1+3w_{0})
             - \frac{9\,\mathrm{KAL}_{0}\,H^{2}}{8\pi G}\right],
  \label{eq:deriv_coeff9}
\end{align}
which is exactly Eq.~(\ref{eq:Friedmann_a}).  The factor 9 arises from
$3 \times |\Pi|/(H^{2}) = 3 \times 3\mathrm{KAL}_{0}/8\pi G$; a factor
of 3 rather than 9 would correspond to including $\Pi$ only once instead
of as the trace of the viscous stress.

Similarly, the conservation law $\nabla_{\!\mu}T^{\mu\nu}=0$ for the effective
fluid $(w_{0},\Pi)$ gives
$\dot{\rho} + 3H(\rho + w_{0}\rho + \Pi) = 0$, hence
$\dot{\rho} + 3H\rho(1+w_{0}) = 9\mathrm{KAL}_{0}H^{3}/8\pi G$,
confirming that Eq.~(\ref{eq:Friedmann_a}) and Eq.~(\ref{eq:continuity})
are mutually consistent.

% ─────────────────────────────────────────────────────────────────────────────
\subsection{Israel--Stewart Viscous Perturbations}
\label{sec:eft:IS}

The Eckart formulation used in \S\ref{sec:eft:PX}--\ref{sec:eft:friedmann}
is a first-order theory that violates causality \citep{HiscockLindblom1985}:
the viscous signal propagates instantaneously.  Israel--Stewart (IS) theory
\citep{IsraelStewart1979} resolves this by promoting $\Pi$ to a dynamical
variable with a finite relaxation time $\tau_{\Pi}$:
%
\begin{equation}
  \tau_{\Pi} H_{0}\,\tilde{h}(a)\,
    \frac{\mathrm{d}\tilde{\Pi}}{\mathrm{d}\ln a}
  + \tilde{\Pi}
  = -\,\mathrm{KAL}_{0}\,\tilde{h}^{2}(a),
  \label{eq:IS_transport}
\end{equation}
%
where $\tilde{h} = H/H_{0}$, $\tilde{\Pi} = \Pi\,8\pi G/H_{0}^{2}$, and
all densities are normalised to $\rho_{\mathrm{crit},0}$.  In the
limit $\tau_{\Pi} \to 0$, Eq.~(\ref{eq:IS_transport}) reduces to the
Eckart result $\tilde{\Pi} = -\mathrm{KAL}_{0}\tilde{h}^{2}$.

\paragraph{Causality constraint.}
The bulk-viscous sound speed in IS theory is \citep{HiscockLindblom1985}
%
\begin{equation}
  c_{s}^{2}
    = \frac{\zeta}{\rho_{\mathrm{DE}}\,\tau_{\Pi}}
    = \frac{\mathrm{KAL}_{0}}{3\,\Omega_{\mathrm{DE}}\,(\tau_{\Pi} H_{0})}
    \leq 1.
  \label{eq:IS_causality}
\end{equation}
%
Setting $c_{s}^{2} = 1$ (the causality boundary) uniquely fixes the
IS relaxation time as an algebraic combination of SSEE constants:
%
\begin{equation}
  \tau_{\Pi} H_{0}
    \;=\; \frac{\mathrm{KAL}_{0}}{3\,\Omega_{\mathrm{DE}}}
    \;=\; \frac{\mathrm{KAL}_{0}\,M_{v}}{3\,T_{r}}
    \;\approx\; 2.191,
  \label{eq:tau_IS}
\end{equation}
%
with no free parameters.

\paragraph{Growth index.}
Coupling Eq.~(\ref{eq:IS_transport}) to the conservation
equation~(\ref{eq:continuity}) and the linear-growth ODE
%
\begin{equation}
  \delta'' + \left[2 + \frac{\mathrm{d}\ln\tilde{h}}{\mathrm{d}\ln a}\right]\delta'
  = \frac{3}{2}\,\frac{\Omega_{m,\mathrm{dyn}}\,a^{-3}}{\tilde{h}^{2}(a)}\,\delta,
  \label{eq:growth_IS}
\end{equation}
%
we solve the system numerically (\texttt{ssee\_is\_growth.py}) with
the IS background $\tilde{h}^{2}(a) = \Omega_{m,\mathrm{dyn}}\,a^{-3}+\rho_{\mathrm{DE}}(a)$.
The growth rate $f \equiv \mathrm{d}\ln\delta/\mathrm{d}\ln a \approx \Omega_{m}(a)^{\gamma}$
is well described by a power law over $0.1 \leq a \leq 1$ with
%
\begin{equation}
  \gamma_{\mathrm{IS}} = 0.657 \pm 0.002
  \quad(\text{essentially independent of }\tau_{\Pi} H_{0}).
  \label{eq:gamma_IS}
\end{equation}
%
This differs by 6\% from the algebraic sovereignty $\gamma = \varphi^{-1} = 0.618$.
Paper~3 \S5.4 reports both values: the algebraic prediction $\gamma_\mathrm{alg}=0.618$
and the IS numerical result $\gamma_\mathrm{IS}=0.657$.
The IS formalism provides the causal theoretical structure; whether $\gamma = \varphi^{-1}$
is the true attractor of the IS boundary conditions remains a target prediction
for forthcoming Stage-IV weak-lensing surveys (LSST, Euclid).
Note that the direct IS growth-factor integral (this Appendix) yields $S_8=0.837$,
while Paper~3's CAMB post-processing method with the same $\gamma_\mathrm{IS}=0.657$
gives $S_8=0.818$; both are consistent with Planck 2018 ($0.832\pm0.013$) within $1\sigma$.

\paragraph{$S_{8}$ constraint.}
Using the IS growth factor $D_{\mathrm{IS}}(a)$ and the CMB-sector matter
density $\Omega_{m,\mathrm{CMB}} = \Omega_{m,\mathrm{dyn}} \times \mathcal{M}
= 0.3199$ (MIRA, Paper~3),
%
\begin{equation}
  S_{8}^{\mathrm{IS}}
    \equiv \sigma_{8,\mathrm{IS}}
           \sqrt{\frac{\Omega_{m,\mathrm{CMB}}}{0.3}}
    \approx 0.837,
  \label{eq:S8_IS}
\end{equation}
%
within $0.4\sigma$ of the Planck-2018 value $0.832 \pm 0.013$~\citep{PlanckVI2020}.
Using the dynamic sector $\Omega_{m,\mathrm{dyn}} = 0.160$ instead gives
$S_{8} \approx 0.59$ ($6\sigma$ tension); the correct physical choice depends
on the formal resolution of the two-sector Friedmann problem (Section~\ref{sec:eft:status}, item 3).

% ─────────────────────────────────────────────────────────────────────────────
\subsection{Physical Regimes}
\label{sec:eft:regimes}

\paragraph{Late universe ($z \lesssim 2$, DESI~DR2).}
At low redshift $H \approx H_{0}$, so the viscous source
$\propto \mathrm{KAL}_{0}H^{3}/8\pi G$ is small compared to
$\rho_{\mathrm{SSEE}}$.  The acceleration equation is dominated by
$\rho_{\mathrm{SSEE}}(1+3w_{0}) < 0$
(since $w_{0} \approx -0.840 < -1/3$), driving $\ddot{a} > 0$
without dark matter.  This regime is the one tested against DESI~DR2
in Paper~2, where the algebraic prediction $(w_{0},w_{a})$ matches
the data at $0.05\sigma$.

\paragraph{Early universe (CMB epoch, $z \sim 10^{3}$).}
At recombination $H \gg H_{0}$, the geometric retention term
$9\,\mathrm{KAL}_{0}H^{3}/8\pi G$ acts as a dissipative source in
Eq.~(\ref{eq:continuity}), injecting energy into the dark-energy fluid
and modifying its effective equation of state.  This structural
viscosity provides the physical mechanism by which an effective
matter density $\Omega_{m,\mathrm{eff}} = 1 - T_{r}/M_{v} \approx 0.160$
yields the correct CMB acoustic scale through the MIRA factor
$\mathcal{M} = 2 - 1/\varphi^{2} \approx 1.9989$ (Paper~3, \S2.3).

% ─────────────────────────────────────────────────────────────────────────────
% ─────────────────────────────────────────────────────────────────────────────
\subsection{Inflationary Embedding: \texorpdfstring{$\alpha$}{α}-Attractor with \texorpdfstring{$\alpha=\varphi^{4}/3$}{α=φ⁴/3}}
\label{sec:eft:inflation}

The SSEE algebraic predictions for the primordial scalar spectral index
$n_s = 1 - \varphi^{-7}$ (Paper~1, \S\ref{sec:register}) and the
tensor-to-scalar ratio $r = \varphi^{-10}$ (Paper~4) belong simultaneously
to a well-defined class of inflationary models: the
\emph{$\alpha$-attractors} of Kallosh \& Linde \citeyearpar{KalloshLinde2013}.

\paragraph{$\alpha$-attractor predictions.}
In the universal leading-order limit, $\alpha$-attractor models predict
\begin{equation}
  n_s = 1 - \frac{2}{N}, \qquad
  r   = \frac{12\alpha}{N^{2}},
  \label{eq:alpha_attractor}
\end{equation}
where $N$ is the number of inflationary e-folds and $\alpha > 0$
parameterises the inverse curvature of the scalar field's K\"{a}hler target
manifold \citep{KalloshLinde2013,Ferrara2013}.
Setting $\alpha = 1$ recovers Starobinsky $R^{2}$ inflation exactly.

\paragraph{SSEE identification.}
Inserting $n_s = 1 - \varphi^{-7}$ into the first relation gives
\begin{equation}
  N = 2/\varphi^{-7} = 2\varphi^{7} \approx 58.07,
  \label{eq:N_ssee}
\end{equation}
which lies squarely in the observationally required window $50\lesssim N\lesssim 60$.
Substituting $N = 2\varphi^{7}$ and $r = \varphi^{-10}$ into the second relation:
\begin{equation}
  \alpha = \frac{r\,N^{2}}{12}
         = \frac{\varphi^{-10}\cdot 4\varphi^{14}}{12}
         = \frac{4\varphi^{4}}{12}
         = \frac{\varphi^{4}}{3}.
  \label{eq:alpha_ssee}
\end{equation}
This result is algebraically exact (verified numerically to machine precision;
see \texttt{ssee\_inflation\_connection.py}).
Using the Fibonacci identity $\varphi^{4} = 3\varphi + 2$:
\begin{equation}
  \alpha = \frac{3\varphi + 2}{3} = \varphi + \frac{2}{3} \approx 2.285.
  \label{eq:alpha_fibonacci}
\end{equation}

\paragraph{Geometric interpretation.}
In the $\alpha$-attractor framework the K\"{a}hler curvature of the
inflaton's target space (a Poincar\'{e} disk of radius $\sqrt{3\alpha}$)
is $\mathcal{R} = -2/(3\alpha)$.
For $\alpha = \varphi^{4}/3$:
\begin{equation}
  \mathcal{R} = -\frac{2}{3 \cdot \varphi^{4}/3} = -\frac{2}{\varphi^{4}}
               = -\varphi^{-4} \approx -0.146.
  \label{eq:kahler_curvature}
\end{equation}
The inflaton therefore lives on a hyperbolic manifold whose curvature is
set by $\varphi$ alone. This gives a \emph{geometric} origin for the
appearance of $\varphi$ in inflationary observables: $\varphi$ enters not
as a numerical coincidence but as the intrinsic length scale of the
K\"{a}hler geometry of the SSEE inflaton sector.

\paragraph{Slow-roll parameters.}
The SSEE predictions $(n_s,\,r) = (1-\varphi^{-7},\,\varphi^{-10})$ fix the
slow-roll parameters uniquely:
\begin{align}
  \varepsilon &= \frac{r}{16} = \frac{\varphi^{-10}}{16}, \label{eq:sr_eps} \\
  \eta &= \frac{n_s - 1 + 6\varepsilon}{2}
        = \frac{-\varphi^{-7} + 3\varphi^{-10}/8}{2}, \label{eq:sr_eta} \\
  \frac{\eta}{\varepsilon} &= 3 - 8\varphi^{3} = -5 - 16\varphi \approx -30.9.
  \label{eq:sr_ratio}
\end{align}
The ratio $\eta/\varepsilon = 3 - 8\varphi^{3}$ is a pure algebraic
identity derived from the SSEE axioms $\{\varphi, \pi\}$; it defines a
one-parameter family of inflationary potentials $V(\phi_{\rm inf})$ whose
single element consistent with $\alpha$-attractors is $\alpha = \varphi^{4}/3$.

\paragraph{Relation to Starobinsky.}
Starobinsky ($\alpha = 1$) reproduces $n_s$ correctly for $N = 2\varphi^{7}$
but gives $r_{\rm Staro} = 3/\varphi^{14} \approx 0.00356$, a factor of
$\varphi^{4}/3\approx 2.28$ below the SSEE prediction $\varphi^{-10}$.
SSEE is therefore \emph{not} Starobinsky inflation; it is a generalised
$\alpha$-attractor that reduces to Starobinsky only in the limit
$\varphi \to 1$ (i.e.\ $\alpha \to 1$, the fixed point of the golden-ratio
recursion $x = 1 + 1/x$).

% ─────────────────────────────────────────────────────────────────────────────
\subsection{Summary and Open Problems}
\label{sec:eft:status}

Table~\ref{tab:eft_summary} summarises the EFT status.
The action~(\ref{eq:eft_action}) with kinetic exponent~(\ref{eq:n_ssee})
and viscosity~(\ref{eq:zeta_ssee}) constitutes a \emph{complete},
zero-free-parameter EFT description of the SSEE background evolution.
Every constant ($n$, $\mathrm{KAL}_{0}$, $w_0$, $w_a$, $\Omega_{\mathrm{DE}}$) is
derived from $\varphi$ and $\pi$ alone; only the overall energy scale $H_{0}$
enters as an observational input.  Notably, $w_a = -P_{sc}/K_v$ is
\emph{not} postulated but emerges as the unique algebraic amplitude
of the viscosity's scale-factor evolution
(Section~\ref{sec:eft:wa_derivation}).

Four aspects require further development:
%
\begin{enumerate}
  \item \textbf{Growth index.}
        The IS-derived growth index $\gamma_{\mathrm{IS}} = 0.657$
        (\S\ref{sec:eft:IS}) differs by 6\% from the algebraic prediction
        $\varphi^{-1} = 0.618$.  A mechanism that would select
        $\gamma = \varphi^{-1}$ from IS boundary conditions is not yet
        identified; this constitutes a near-prediction testable by
        LSST/Euclid Stage-IV weak lensing.

  \item \textbf{Full perturbation spectrum.}
        The IS growth ODE~(\ref{eq:growth_IS}) governs sub-Hubble dark-matter
        perturbations; the matching to CMB scalar transfer functions and
        the tensor-mode propagation in the IS background require a complete
        Boltzmann implementation.

  \item \textbf{Radiative stability.}
        The EFT cutoff scale $\Lambda_{\mathrm{UV}}$ and the quantum stability
        of the $n < 0$ power-law sector are not yet established.

  \item \textbf{Inflation--dark energy unification.}
        The $\alpha$-attractor sector (\S\ref{sec:eft:inflation}) and the
        k-essence dark-energy sector (\S\ref{sec:eft:PX}) both originate
        from $\varphi$, suggesting a single unified action.
        Constructing the quintessential inflation potential $V(\phi)$
        that reduces to $P(X) = \mathcal{A}X^n$ at late times
        and to an $\alpha=\varphi^4/3$ attractor at early times
        is identified as high-priority future work.
\end{enumerate}

\begin{table}[h]
\centering
\caption{EFT status summary for SSEE-V3.6.}
\label{tab:eft_summary}
\begin{tabular}{lll}
\hline\hline
Element & Status & Derivation \\
\hline
Action form $S$ & \checkmark Complete & Standard GR + k-essence + Eckart \\
$u_{\mu}(\phi)$ coupling & \checkmark Complete & $u_{\mu} = -\partial_{\mu}\phi/\sqrt{2X}$ \\
Kinetic exponent $n$ & \checkmark Algebraic & $n=(T_{r}-M_{v})/2T_{r}$ from $\varphi,\pi$ \\
Equation of state $w_{0}$ & \checkmark Algebraic & $w_{0}=1/(2n-1)=-T_{r}/M_{v}$ \\
Viscosity coupling $\zeta$ & \checkmark Algebraic & $\zeta=\mathrm{KAL}_{0}H/8\pi G$ \\
Friedmann eqs.\ I--III & \checkmark Consistent & Factor 9 verified in \S\ref{sec:eft:friedmann} \\
$w_{a}$ from EFT & \checkmark Algebraic & $\zeta(a)\!\propto\!(1{-}a)\rho_\mathrm{DE}/H$; see \S\ref{sec:eft:wa_derivation} \\
IS relaxation time $\tau_\Pi$ & \checkmark Algebraic & $\tau_\Pi H_0 = \mathrm{KAL}_0/(3\Omega_\mathrm{DE}) \approx 2.191$ \\
IS causality ($c_s^2 \leq 1$) & \checkmark Verified & Eq.~(\ref{eq:IS_causality}); $c_s^2 = 1$ at boundary \\
Growth index $\gamma_\mathrm{IS}$ & \checkmark Numerical & $0.657 \pm 0.002$ (6\% above $\varphi^{-1}=0.618$) \\
$S_8$ (CMB sector) & \checkmark Numerical & $0.837$, within $0.4\sigma$ of Planck \\
$\alpha$-attractor ($\alpha=\varphi^4/3$) & \checkmark Algebraic & Eq.~(\ref{eq:alpha_ssee}); exact, $N=2\varphi^7$ \\
K\"{a}hler curvature $\mathcal{R}=-\varphi^{-4}$ & \checkmark Algebraic & Eq.~(\ref{eq:kahler_curvature}) \\
Full perturbation spectrum & $\circ$ Open & Off-diagonal CMB cov.\ + Stage-IV WL \\
Inflation--DE unification & $\circ$ Open & Quintessential inflation potential \\
Radiative stability & $\circ$ Open & Future work \\
\hline
\end{tabular}
\end{table}


%--------------------------------------------------------------------
\section{Dark Energy Evolution}
\label{sec:de}
%--------------------------------------------------------------------

\SSEE\ models dynamical dark energy via the Chevallier–Polarski–Linder (CPL) form
$w(a)=\wz+\wa(1-a)$. Motivated by Eq.~\eqref{eq:action}, the expansion is subject to
an effective geometric pressure:
\begin{equation}
  H^2 = \frac{8\pi G}{3}\!\left(\rho_{\rm bar}+\rho_{\rm DE}\right),
  \qquad
  \rho_{\rm DE} = \rho_{{\rm crit},0}\!\left(\frac{T_r}{M_v}\right).
\end{equation}
The baseline state $\wz$ is set by the static geometric ratio $T_r/M_v$, and the evolution
$\wa$ encodes the dynamical departure from saturation driven by $P/K_v$:
\begin{equation}
  \wz = -\frac{T_r}{M_v} \approx -0.840,
  \qquad
  \wa = -\frac{P}{K_v}  \approx -0.670.
  \label{eq:ssee_pars}
\end{equation}

\noindent\textbf{Explicit reduction to the two axioms.}
Substituting the algebraic definitions
$T_r = \tfrac{3(3\varphi+\pi)}{2}$ (TRIAL),
$M_v = 3(\varphi+\pi)$ (MIKAEL\_V),
$P_{sc} = 2\varphi+\pi$ (PYROS), and
$K_v = 2(\varphi+\pi)$ (KRYSTOS), both
dark-energy observables collapse to closed forms in $\varphi$ and $\pi$ alone:
\begin{equation}
  \boxed{w_0 = -\frac{3\varphi+\pi}{2(\varphi+\pi)} = -0.8399},
  \qquad
  \boxed{w_a = -\frac{2\varphi+\pi}{2(\varphi+\pi)} = -0.6699}.
  \label{eq:w0wa_explicit}
\end{equation}
No fitting procedure is involved: these values are uniquely fixed by the geometry of $\varphi$ and $\pi$.
The numerators $3\varphi+\pi$ and $2\varphi+\pi$ reflect the three- and two-fold
contributions of $\varphi$ to the saturation horizon, while the common denominator
$2(\varphi+\pi)$ is twice the Stability Metric $\Omega=\varphi+\pi$.

\medskip
\noindent\textbf{Asymptotic saturation.}\ Although $\wz+\wa=-1.510$ would ordinarily
imply a phantom Big Rip, \SSEE\ adopts CPL only as a local observational parameterisation.
In the far future ($a\to\infty$) the energy density saturates at the finite $T_r/M_v$
boundary of Eq.~\eqref{eq:action}, ensuring universal stability.

DESI~DR2 \citep{AbdulKarim2025} analyses indicate an evolving preference for dynamical
dark energy. The \SSEE\ algebraic point $(\wz,\wa)=(-0.840,-0.670)$ falls within the
68\% confidence contours of the combined BAO+CMB+DESY5 datasets
(Mahalanobis $\chi^2_{2D}=0.080$, i.e.\ $0.05\sigma$; quantified in Paper~2
\citep{SSEE_paperII}).

%--------------------------------------------------------------------
\section{Galaxy Cluster Dynamics}
\label{sec:clusters}
%--------------------------------------------------------------------

\subsection{The Cluster Mass Discrepancy and IGIMF}

In rich galaxy clusters, inferred dynamical mass typically exceeds the visible baryonic
mass by a factor of $\sim$5--6. Recent literature incorporating a variable IGIMF
\citep{Zhang2026} demonstrates that baryonic masses have been systematically
underestimated: heavy stellar remnants raise the effective baryonic baseline to at least
$88^{+5}_{-4}\%$ of the MOND dynamical mass. In \SSEE, this updated IGIMF baseline
multiplied by $\KALz \approx 5.5214$ natively closes the remaining discrepancy without
cold dark matter.

\subsection{Structural Viscosity and the Neutrino Bridge}

At the cluster scale, the total effective dynamical mass $M_{\rm dyn}$ within radius $R$ is:
\begin{equation}
  M_{\rm dyn}(R) = \Migimf(R)\cdot\KALz\cdot\left[1+f_\nu\right],
  \label{eq:mdyn}
\end{equation}
where $f_\nu$ quantifies the fractional mass contribution of trapped relic neutrinos.
Using $\sum m_\nu=0.085$\,eV — robustly permitted within \SSEE's dynamic dark-energy
regime — the cosmological neutrino density is $\rho_\nu\approx28.6$\,eV\,cm$^{-3}$.
For a cluster with escape velocity $v_{\rm esc}\approx2000$\,km\,s$^{-1}$, integration
of the Fermi–Dirac distribution \citep{Tremaine1979} yields a trapped fraction
$f_\nu\approx0.018$--$0.022$. This structurally consistent $\sim$2\% correction provides
the final macroscopic bridge.

\subsection{Quantitative Cluster Analysis}

Applying Eq.~\eqref{eq:mdyn} to four well-documented clusters demonstrates predictive
accuracy against observed dynamical masses (Table~\ref{tab:clusters}).

\begin{table}[H]
\centering
\caption{Cluster mass predictions. All masses within $R_{200}$ in units of
  $10^{14}\,M_\odot$. $M_{\SSEE} = \Migimf \times 5.5214 \times 1.02$.
  IGIMF baselines and observed masses from \citet{Zhang2026}; Bullet lensing
  mass from \citealt{Clowe2006}.}
\label{tab:clusters}
\small
\begin{tabular}{lrrr}
\toprule
Cluster & $\Migimf$ [$10^{14}\,M_\odot$] & $M_{\SSEE}$ [$10^{14}\,M_\odot$] & $\Mobs$ [$10^{14}\,M_\odot$] \\
\midrule
Coma          & $1.8\pm0.2$ & $10.1\pm1.1$ & $9.8\pm1.0$  \\
A2029         & $2.2\pm0.2$ & $12.4\pm1.1$ & $12.0\pm1.2$ \\
A478          & $1.5\pm0.2$ & $8.4\pm1.1$  & $8.0\pm1.0$  \\
Bullet (main) & $1.2\pm0.2$ & $6.7\pm1.1$  & $6.5\pm1.0$  \\
\bottomrule
\end{tabular}
\end{table}

All four predictions agree with observed masses at the $1\sigma$ level ($\chi^2_r=0.122$;
full analysis in Paper~2 \citep{SSEE_paperII}).

\subsection{Gravitational Lensing and the Newtonian Limit}
\label{sec:lensing}

The Bullet Cluster (1E~0657-56) is frequently cited as proof of particulate dark matter
due to the spatial offset between collisional X-ray gas and lensing-mass peaks
\citep{Clowe2006}. \SSEE\ reproduces this topological offset through macroscopic
structural viscosity without collisionless particles.

The structural viscosity acts as an environment-dependent scalar field. Defining
$x = |\nabla\Phi|/a_0$, \SSEE\ adopts an interpolation function analogous to the
canonical MOND simple interpolating function $\mu(x)$:
\begin{equation}
  \mu_{\SSEE}(x) = \frac{x+1}{x+\KALz}
  \implies
  \mathrm{KAL}(x) = \frac{1}{\mu_{\SSEE}(x)} = \frac{x+\KALz}{x+1}.
  \label{eq:KAL}
\end{equation}
For $x\gg1$ (Newtonian limit), $\mathrm{KAL}(x)\to1$. For $x\ll1$ (deep MONDian limit),
$\mathrm{KAL}(x)\to\KALz$. During a high-velocity cluster collision, the stripped X-ray
gas flattens its local potential gradient ($x\ll1$), activating the full $\KALz$
amplification; collisionless galaxies retain compact potential wells ($x\ge1$), coupling
non-linearly.

The effective lensing convergence is:
\begin{equation}
  \kappa(\vec{\theta})
  = \frac{\Sigma_{\SSEE}(\vec{\theta})}{\Sigma_{\rm crit}},
  \quad
  \Sigma_{\SSEE}(\vec{\theta})
  = \int \rho_{\rm bar}^{\rm IGIMF}(\vec{\theta},\ell)\cdot\mathrm{KAL}(x)\,d\ell,
  \quad
  \Sigma_{\rm crit} = \frac{c^2}{4\pi G}\frac{D_S}{D_L D_{LS}}.
\end{equation}
This allows lensing peaks to track galaxies, consistent with QUMOND simulations
\citep{Angus2007}.

%--------------------------------------------------------------------
\section*{Data Availability}
%--------------------------------------------------------------------

All observational data used in this work are publicly available.
DESI~DR2 BAO measurements are from \citet{AbdulKarim2025}
(\url{https://arxiv.org/abs/2503.14738}).
Galaxy-cluster dynamical masses are from \citet{Zhang2026}
(\url{https://arxiv.org/abs/2602.06082}).
Analysis scripts and data files are available at
\url{https://github.com/mikealmeida1721/SSEE}.

%--------------------------------------------------------------------
\section*{Acknowledgements}
%--------------------------------------------------------------------

The author thanks the DESI Collaboration for public DR2 data and
P.~Kroupa and collaborators for the IGIMF cluster re-analysis.

%--------------------------------------------------------------------
\bibliographystyle{plainnat}
\begin{thebibliography}{99}

\bibitem[Abdul-Karim et al.(2025)]{AbdulKarim2025}
Abdul-Karim, M., et al.\ (DESI Collaboration) (2025).
DESI DR2 Results II: Measurements of Baryon Acoustic Oscillations and Cosmological Constraints.
\textit{Phys.\ Rev.\ D} \textbf{112}, 083515.
arXiv:2503.14738.

\bibitem[Almeida Vallejo(2026)]{SSEE_paperII}
Almeida Vallejo, M.~E.\ (2026).
Bayesian MCMC Validation of the Structural Self-Energy Expansion Framework (SSEE-V3.6).
\textit{Paper~2, this series}.

\bibitem[Angus et al.(2007)]{Angus2007}
Angus, G.~W., et al.\ (2007).
On the proof of dark matter, the law of gravity and the mass of neutrinos.
\textit{MNRAS} \textbf{374}, 1248.

\bibitem[Clowe et al.(2006)]{Clowe2006}
Clowe, D., et al.\ (2006).
A direct empirical proof of the existence of dark matter.
\textit{ApJL} \textbf{648}, L109.

\bibitem[Tremaine \& Gunn(1979)]{Tremaine1979}
Tremaine, S., \& Gunn, J.~E.\ (1979).
Dynamical role of light neutral leptons in cosmology.
\textit{Phys.\ Rev.\ Lett.} \textbf{42}, 407.

\bibitem[Zhang et al.(2026)]{Zhang2026}
Zhang, D., Zonoozi, A.~H., \& Kroupa, P.\ (2026).
Revisiting the missing mass problem in MOND for nearby galaxy clusters.
\textit{Phys.\ Rev.\ D} \textbf{113}, 043027.
arXiv:2602.06082.

\bibitem[Eckart(1940)]{Eckart1940}
Eckart, C.\ (1940).
The thermodynamics of irreversible processes. III. Relativistic theory of the simple fluid.
\textit{Phys.\ Rev.} \textbf{58}, 919.

\bibitem[Weinberg(1971)]{Weinberg1971}
Weinberg, S.\ (1971).
Entropy generation and the survival of protogalaxies in an expanding universe.
\textit{ApJ} \textbf{168}, 175.

\bibitem[Israel \& Stewart(1979)]{IsraelStewart1979}
Israel, W., \& Stewart, J.~M.\ (1979).
Transient relativistic thermodynamics and kinetic theory.
\textit{Ann.\ Phys.} \textbf{118}, 341.

\bibitem[Brevik et al.(2005)]{Brevik2005}
Brevik, I., Nojiri, S., Odintsov, S.~D., \& Vanzo, L.\ (2005).
Entropy and universality of the Cardy-Verlinde formula in dark energy universe.
\textit{Phys.\ Rev.\ D} \textbf{70}, 043520.

\bibitem[Hiscock \& Lindblom(1985)]{HiscockLindblom1985}
Hiscock, W.~A., \& Lindblom, L.\ (1985).
Generic instabilities in first-order dissipative relativistic fluid theories.
\textit{Phys.\ Rev.\ D} \textbf{31}, 725.

\bibitem[Planck Collaboration VI(2020)]{PlanckVI2020}
Planck Collaboration (Aghanim, N., et al.)\ (2020).
Planck 2018 results VI: Cosmological parameters.
\textit{A\&A} \textbf{641}, A6.
arXiv:1807.06209.

\bibitem[Kallosh \& Linde(2013)]{KalloshLinde2013}
Kallosh, R., \& Linde, A.\ (2013).
Universality class in conformal inflation.
\textit{JCAP} \textbf{07}, 002.
arXiv:1306.5220.

\bibitem[Ferrara et al.(2013)]{Ferrara2013}
Ferrara, S., Kallosh, R., Linde, A., \& Porrati, M.\ (2013).
Minimal supergravity models of inflation.
\textit{Phys.\ Rev.\ D} \textbf{88}, 085038.
arXiv:1307.7696.

\end{thebibliography}

%=====================================================================
\section{Domain of Validity}
\label{sec:domain}
%=====================================================================

A persistent risk in multi-sector theoretical models is that statements of
validity or falsification are conflated across domains that probe different
physical regimes. This section explicitly delineates the domain of validity
of each SSEE sector and the dataset used to validate it. This delineation
applies uniformly to Papers~1--4.

\begin{description}

  \item[\textbf{Dynamic sector} (low redshift, $z\lesssim2.3$).]
  Governed by $\Omega_{m,\rm dyn}=0.1601$, $w_0=-0.8399$, $w_a=-0.6699$.
  Validated against: DESI~DR2 BAO (13 measurements, $0.295\leq z\leq2.33$),
  galaxy-cluster mass data (IGIMF baseline, Paper~2).
  \textit{Result}: $\Delta\mathrm{BIC}=-5.55$ (dynamic sector, Paper~2),
  $0.05\sigma$ from DESI~DR2 CPL posterior in the $w_0$--$w_a$ plane.
  \textit{Known limitation}: under standard Friedmann background,
  $\chi^2_r(H(z))=1.861$ (Cosmic Chronometers); this is a direct consequence of
  the enlarged $r_d=175.6$~Mpc and is documented, not hidden (Paper~2~\S5.3).

  \item[\textbf{Observational sector} (high redshift, CMB).]
  Governed by $\Omega_{m,\rm CMB}=\Omega_{m,\rm dyn}\times\mathrm{MIRA}=0.3199$,
  where MIRA~$=(3\varphi+\pi)/4=1.99892$ is algebraically pre-committed
  (January 2026, Zenodo~\texttt{10.5281/zenodo.19679049}).
  Validated against: Planck PR4 TT+TE+EE+lensing via CAMB; \texttt{plik\_lite}
  TTTEEE likelihood via Cobaya (Paper~3).
  \textit{Result}: $\chi^2_r=1.062$ (TT, 1971 multipoles);
  $\Delta\mathrm{BIC}=-40.3$ at $H_0=67.075$~km\,s$^{-1}$\,Mpc$^{-1}$ ($k=1$
  vs $k=6$ for $\Lambda$CDM).
  \textit{Claim}: SSEE is \emph{not falsified} by the Planck PR4 CMB spectra.
  This claim is sector-specific: it refers to the observational sector
  (with MIRA applied), not the dynamic sector directly.

  \item[\textbf{Algebraic sector} (purely geometric).]
  All constants fixed from $\{\varphi,\pi\}$ with zero data input.
  Cross-validated against: Planck 2018 ($n_s$, $\Omega_b h^2$, $H_0$),
  LiteBIRD forecast ($r$), BBN ($\sum m_\nu$), Paper~4.
  \textit{Result}: all Type~A predictions in Appendix~B,
  Table~\ref{tab:derivation_chain} are within $2\sigma$ of published
  observational values (see Paper~4 for full compilation).
  \textit{Claim}: the algebraic sector is \emph{not falsified} by any
  currently available dataset. A future falsification would require
  $n_s\neq1-\varphi^{-7}$ at $>3\sigma$ from Planck~PR4,
  or $r\neq\varphi^{-10}$ from LiteBIRD.

\end{description}

\paragraph{On cross-sector BIC comparisons.}
The $\Delta\mathrm{BIC}=-5.55$ (Paper~2, dynamic sector) and
$\Delta\mathrm{BIC}=-40.3$ (Paper~3, observational sector via
\texttt{plik\_lite}) are \emph{not directly comparable}: they operate on
different datasets ($N=16$ vs $N=6413$ data points), under different
background assumptions (Friedmann vs MIRA-corrected CMB), and with
different parameter counts ($k=1$ vs $k=1$, but $\Lambda$CDM has $k=3$
and $k=6$ respectively). Both values are internally valid within their
respective domains; neither generalises to the other.

%=====================================================================
\appendix
%=====================================================================

\section{Master Derivation Chain: $\varphi,\pi \to$ Observables}
\label{app:derivation_chain}

Table~\ref{tab:derivation_chain} provides the complete, unambiguous
derivation chain from the two axioms ($\varphi$, $\pi$) to every
cosmological observable used in Papers~1--4.
This table is the authoritative cross-reference for all four manuscripts.
Each intermediate constant is labelled as either a \textit{primary register}
(derived in one algebraic step from $\varphi$ and $\pi$) or a
\textit{secondary register} (derived from primary registers).

\begin{table}[H]
\centering
\caption{Master derivation chain for the \\SSEE\\ framework.
  Result types: \textbf{A} = pure algebraic (no data input),
  \textbf{M} = MCMC-validated (data-constrained posterior),
  \textbf{P} = prior-dependent (requires external dataset anchor).
  $\varphi = (1+\sqrt{5})/2 \approx 1.61803$;
  $\pi \approx 3.14159$.}
\label{tab:derivation_chain}
\small
\begin{tabular}{llllc}
\toprule
Symbol & Definition & Value & Observable link & Type \\
\midrule
\multicolumn{5}{l}{\textit{Level 0 — Axioms}} \\
$\varphi$ & $(1+\sqrt{5})/2$ & $1.618034$ & All & A \\
$\pi$     & Archimedes constant & $3.141593$ & All & A \\
\midrule
\multicolumn{5}{l}{\textit{Level 1 — Primary Registers}} \\
$\Omega_{\rm DNAV}$ & $\varphi + \pi$ & $4.759627$ & Stability metric & A \\
$\beta$   & $(\varphi+\pi)/2$ & $2.379814$ & Base coupling & A \\
$P_{sc}$  & $\Omega_{\rm DNAV}+\varphi$ & $6.377661$ & Dyn.\ evolution & A \\
AURA      & $\varphi+\beta$ & $3.997838$ & Resonance limit & A \\
\midrule
\multicolumn{5}{l}{\textit{Level 2 — Secondary Registers}} \\
$\mathrm{KAL}_0$ & $\beta+\pi$ & $5.521406$ & Viscosity; $M_{\rm cluster}$ & A \\
$K_v$     & $\varphi+\pi+\Omega_{\rm DNAV}$ & $9.519254$ & Struct.\ constraint & A \\
$T_r$     & $3(\varphi+\beta)$ & $11.993520$ & Saturation horizon & A \\
MIRA      & $\mathrm{AURA}/2 = (\varphi+\beta)/2$ & $1.998924$ & CMB--dyn.\ bridge & A$^*$ \\
$M_v$     & $\varphi+\pi+K_v$ & $14.278880$ & Max.\ dim.\ invariant & A \\
\midrule
\multicolumn{5}{l}{\textit{Level 3 — Cosmological Observables}} \\
$w_0$     & $-T_r/M_v$ & $-0.8399$ & DESI DR2 & A \\
$w_a$     & $-P_{sc}/K_v$ & $-0.6699$ & DESI DR2 & A \\
$H_0^{\rm alg}$ & $3(\varphi+\pi)^2$ & $67.962$ km/s/Mpc & Planck 2018 & A \\
$H_0^{\rm MCMC}$ & emcee posterior & $66.75\pm0.44$ & DESI+Planck & M \\
$\Omega_m^{\rm CMB}$ & $(\pi-\varphi)/(\pi+\varphi)$ & $0.3201$ & Planck 2018 & A \\
$n_s$     & $1-\varphi^{-7}$ & $0.96556$ & Planck 2018 & A \\
$\Omega_b h^2$ & $3(\pi-\varphi)/200$ & $0.02285$ & Planck 2018 & A \\
$\Omega_c h^2$ & $\mathrm{KAL}_0\times\Omega_b h^2\times n_s$ & $0.11926$ & Planck 2018 & A \\
$r$       & $\varphi^{-10}$ & $0.00813$ & LiteBIRD 2032 & A \\
$\sum m_\nu$ & (Acoustic Residue) & $0.084$ eV & Planck+BBN & A \\
\bottomrule
\end{tabular}
\\[4pt]
\small $^*$ MIRA $= (\varphi+\beta)/2 = (3\varphi+\pi)/4 = 1.998924\ldots$
is the canonical value used in all four manuscripts.
It serves a dual role: (i) as a \textit{derived constant} (algebraically
fixed from $\{\varphi,\pi\}$ via $\beta=(\varphi+\pi)/2$, committed January 2026
before any CMB analysis; see Paper~3~\S\ref{sec:mira_timeline}), and (ii) as a
\textit{sector-mapping operator} projecting $\Omega_{m,\rm dyn}=0.1601$
onto $\Omega_{m,\rm CMB}=0.3199$ via
$\Omega_{m,\rm CMB} = \Omega_{m,\rm dyn}\times\mathrm{MIRA} = 0.3199$.
Note: in early internal documents MIRA was loosely approximated as
$\varphi\approx1.618$; the canonical value is $1.998924$ throughout all four papers.
Similarly, $\mathrm{KAL}_0$ is both a derived constant (Level~2) and
a dynamical viscosity operator scaling cluster masses.
\end{table}

%---------------------------------------------------------------------
\section{Unified Symbol Glossary}
\label{app:notation}

To avoid ambiguity across Papers~1--4, Table~\ref{tab:notation}
provides the canonical definition of every non-standard symbol used
in this series.

\begin{table}[H]
\centering
\caption{Unified symbol glossary for the \\SSEE\\ four-paper series.
  All symbols are dimensionless unless otherwise stated.}
\label{tab:notation}
\small
\begin{tabular}{lll}
\toprule
Symbol & Full name & First defined \\
\midrule
$\varphi$ & Golden ratio $(1+\sqrt{5})/2$ & Paper 1 §1 \\
$\Omega_{\rm DNAV}$ & Dark Navigation scalar ($\varphi+\pi$) & Paper 1 §1 \\
$\beta$ & Base coupling scalar & Paper 1 §1 \\
AURA & Acoustic Unified Resonance Amplitude ($2\varphi$) & Paper 1 §2 \\
MIRA & Mapped Inference Resonance Amplitude (AURA/2 $=\varphi$) & Paper 3 §2.4 \\
$\mathrm{KAL}_0$ & Asymptotic structural viscosity factor ($\beta+\pi$) & Paper 1 §1 \\
$P_{sc}$ & Dynamical evolution scalar ($\Omega_{\rm DNAV}+\varphi$) & Paper 1 §1 \\
$K_v$ & Structural constraint invariant & Paper 1 §1 \\
$T_r$ & 3D saturation horizon ($3(\varphi+\beta)$) & Paper 1 §1 \\
$M_v$ & Maximal dimensional invariant ($\varphi+\pi+K_v$) & Paper 1 §1 \\
PYROS & Radiation sector coupling ($\pi^2/6$) & Paper 4 §3 \\
ÎCËBÊRG & $M_v$ (Nine Sovereignties sum) & Paper 4 §4 \\
$\Omega_{m,\rm dyn}$ & Dynamic matter density (SSEE sector) & Paper 2 §2 \\
$\Omega_{m,\rm CMB}$ & CMB-observed matter density ($2\,\mathrm{MIRA}\,\Omega_{m,\rm dyn}$) & Paper 3 §2.4 \\
$\gamma_{\rm IS}$ & Israel--Stewart growth index ($0.657\pm0.002$) & Paper 1 App.A \\
$\tau_\Pi H_0$ & IS relaxation time (dimensionless, $\approx2.191$) & Paper 1 App.A \\
\bottomrule
\end{tabular}
\end{table}

%---------------------------------------------------------------------
\section{Result-Type Classification Protocol}
\label{app:result_types}

A recurring risk in multi-paper series is that results of different
epistemic status are compared without adequate labelling.
To ensure methodological transparency across the four manuscripts,
all numerical results in this series are classified into one of three
categories (see also Table~\ref{tab:derivation_chain}, column ``Type''):

\begin{itemize}
  \item \textbf{Type A — Algebraic:} The value follows solely from
  $\varphi$ and $\pi$ through the Level~1--3 chain in
  Table~\ref{tab:derivation_chain}. No observational data enters the
  computation. These results are \emph{prior-free predictions}.

  \item \textbf{Type M — MCMC-validated:} The value is the median of
  the posterior distribution from the 100-walker \texttt{emcee}
  pipeline described in Paper~2, using DESI~DR2 BAO, Planck~2018
  priors, and galaxy-cluster mass data. These results reflect the
  empirical range accessible within the SSEE parameter family.

  \item \textbf{Type P — Prior-dependent:} The value depends on an
  external dataset anchor (e.g., Planck TT+TE+EE+lensing as a
  likelihood) that is not internally derived by the SSEE framework.
  Type~P results quantify the consistency of SSEE with standard
  cosmological datasets; they are not Type~A predictions.
\end{itemize}

The distinction is particularly important for the $\Delta\mathrm{BIC}$
comparisons in Papers~2 and~3.
In Paper~2, the $\Delta\mathrm{BIC}=-5.55$ (dynamic sector) is a
Type~M result: it reflects the MCMC posterior relative to $\Lambda$CDM
under identical priors.
In Paper~3, the $\Delta\mathrm{BIC}=-40.3$ (plik\_lite likelihood)
is a Type~P result: it depends on the \texttt{plik\_lite} likelihood
as an external CMB anchor.
These two BIC values are \emph{not directly comparable} because they
operate in different model spaces; the distinction is disclosed
explicitly in Paper~3~§4.4.

\end{document}
